% Problems, Chapter 5: Vector Spaces
% Arfken, Weber & Harris, Mathematical Methods for Physicists, 7th ed.
% Requires: amsmath, amssymb, tikz

\section*{Problems}

\Needspace{9\baselineskip}
\subsection*{5.1 Vectors in Function Spaces}

\Needspace{5\baselineskip}
\par\bigskip
\noindent\textbf{5.1.1}\quad A function $f(x)$ is expanded in a series of
orthonormal functions, \erratum{with convergence in the Hilbert-space norm},
\[
f(x)=\sum_{n=0}^{\infty}a_{n}\varphi_{n}(x).
\]
Show that the series expansion is unique for a given set of $\varphi_{n}(x)$.
The functions $\varphi_{n}(x)$ are being taken here as the basis vectors in an
infinite-dimensional Hilbert space.
\solution
Write $S_N=\sum_{n=0}^N a_n\varphi_n$; the expansion means that
$S_N\to f$ in the Hilbert-space norm.
\erratum{The textbook does not say in what sense the series converges,
and uniqueness depends on it: the argument below uses norm convergence,
whereas for convergence almost everywhere uniqueness fails (Menshov
constructed a trigonometric series, with coefficients not all zero,
that converges to zero almost everywhere).}
For a fixed $k$ and $N\ge k$,
\[
\langle\varphi_k|S_N\rangle
=\sum_{n=0}^N a_n\langle\varphi_k|\varphi_n\rangle
=\sum_{n=0}^N a_n\delta_{kn}=a_k.
\]
The scalar product is continuous: by linearity in the second argument and
Schwarz's inequality, for every $N\ge k$,
\[
|\langle\varphi_k|f\rangle-a_k|
=|\langle\varphi_k|f-S_N\rangle|
\le\|\varphi_k\|\,\|f-S_N\|=\|f-S_N\|\longrightarrow0.
\]
The left side does not depend on $N$, so it is zero. Each coefficient is
therefore determined by $f$ alone:
\[
\boxed{a_k=\langle\varphi_k|f\rangle.}
\]
Any other convergent expansion in the same orthonormal set has the same
scalar products and hence the same coefficients. Equality of functions in
this Hilbert space means equality almost everywhere; changes at isolated
points do not produce a different Hilbert-space vector.

\Needspace{5\baselineskip}
\par\bigskip
\noindent\textbf{5.1.2}\quad A function $f(x)$ is represented by a finite set
of basis functions $\varphi_{i}(x)$,
\[
f(x)=\sum_{i=1}^{N}c_{i}\varphi_{i}(x).
\]
Show that the components $c_{i}$ are unique, that no different set $c'_{i}$
exists.

\emph{Note.}~Your basis functions are automatically linearly independent. They
are not necessarily orthogonal.
\solution
Suppose also $f=\sum_{i=1}^N c'_i\varphi_i$. Subtract the two finite sums:
\[
0=\sum_{i=1}^N(c_i-c'_i)\varphi_i.
\]
Linear independence means precisely that a zero linear combination of the
$\varphi_i$ has every coefficient zero. Consequently $c_i-c'_i=0$ for
$i=1,\ldots,N$, proving uniqueness without any orthogonality assumption.

As a scalar-product check, let $G_{ij}=\langle\varphi_i|\varphi_j\rangle$.
For any nonzero column $d$, linear independence makes $\sum_id_i\varphi_i$
a nonzero function, so
\[
d^\dagger Gd=\sum_{i,j}d_i^*\langle\varphi_i|\varphi_j\rangle d_j
=\left\|\sum_i d_i\varphi_i\right\|^2>0.
\]
Hence $Gd\ne0$ whenever $d\ne0$, i.e., $G$ is invertible. Taking the
scalar product of the expansion of $f$ with $\varphi_i$ gives the equations
$\sum_jG_{ij}c_j=\langle\varphi_i|f\rangle$, which therefore also have
exactly one solution.

\Needspace{5\baselineskip}
\par\bigskip
\noindent\textbf{5.1.3}\quad A \erratum{square-integrable} function $f(x)$ is approximated by a power
series $\sum_{i=0}^{n-1}c_{i}x^{i}$ over the interval $[0,1]$. Show that
minimizing the mean square error leads to a set of linear equations
\[
\mathsf{A}\mathbf{c}=\mathbf{b},
\]
where
\[
A_{ij}=\int_{0}^{1}x^{i+j}\,dx=\frac{1}{i+j+1},\quad i,j=0,1,2,\ldots,n-1
\]
and
\[
b_{i}=\int_{0}^{1}x^{i}f(x)\,dx,\quad i=0,1,2,\ldots,n-1.
\]
\emph{Note.}~The $A_{ij}$ are the elements of the Hilbert matrix of order $n$.
The determinant of this Hilbert matrix is a rapidly decreasing function of
$n$. For $n=5$, $\det\mathsf{A}=3.7\times10^{-12}$ and the set of equations
$\mathsf{A}\mathbf{c}=\mathbf{b}$ is becoming ill-conditioned and unstable.
\solution
\erratum{The added square integrability is needed: otherwise the mean
square error is infinite for every choice of the $c_i$, because
$|f|^2\le2|f-p|^2+2|p|^2$ and the polynomial $p=\sum_ic_ix^i$ is bounded
on $[0,1]$.
It also makes each $b_i$ finite, by Schwarz's inequality
$|b_i|\le(\int_0^1x^{2i}dx)^{1/2}(\int_0^1|f|^2dx)^{1/2}$.}
For real $f$ and real coefficients, put $p(x)=\sum_{j=0}^{n-1}c_jx^j$.
Expanding the error before differentiating gives
\begin{align*}
E(c)&=\int_0^1(f-p)^2\,dx\\
&=\int_0^1 f^2\,dx
 -2\sum_jc_j\int_0^1x^jf\,dx
 +\sum_{j,k}c_jc_k\int_0^1x^{j+k}\,dx.
\end{align*}
In the notation of the statement,
$E(c)=\int_0^1f^2dx-2\sum_jb_jc_j+\sum_{j,k}A_{jk}c_jc_k$. The last matrix
is symmetric, $A_{jk}=A_{kj}$, so
$\partial(\sum_{j,k}A_{jk}c_jc_k)/\partial c_i
=\sum_kA_{ik}c_k+\sum_jA_{ji}c_j=2\sum_jA_{ij}c_j$, and
\[
\frac{\partial E}{\partial c_i}
=-2b_i+2\sum_{j=0}^{n-1}A_{ij}c_j=0,
\qquad
A_{ij}=\left[\frac{x^{i+j+1}}{i+j+1}\right]_0^1
=\frac1{i+j+1}.
\]
Hence the normal equations are $\boxed{Ac=b}$. They determine the global
minimum, not merely a stationary point: for any nonzero vector $v$,
\[
v^TAv=\int_0^1\left(\sum_{j=0}^{n-1}v_jx^j\right)^2dx>0.
\]
A nonzero polynomial cannot vanish on an interval, so the Hessian $2A$ is
positive definite. In particular $Av\ne0$ for $v\ne0$, so $A$ is invertible
and the normal equations have exactly one solution $c_*$. To complete the
square, put $c=c_*+d$ in $E(c)=\int_0^1f^2dx-2b^Tc+c^TAc$ and use
$b=Ac_*$, $A^T=A$:
\begin{align*}
E(c_*+d)&=\int_0^1f^2dx-2c_*^TA(c_*+d)+(c_*+d)^TA(c_*+d)\\
&=\int_0^1f^2dx-c_*^TAc_*+d^TAd=E(c_*)+d^TAd.
\end{align*}
Since $d^TAd>0$ for $d\ne0$, $c_*$ is the unique minimizer.

For complex $f$ and complex $c$, minimize $E(c)=\int_0^1|f-p|^2dx$. The
$x^j$ are real, so $\int_0^1f^*p\,dx=b^\dagger c$ and
$\int_0^1fp^*dx=c^\dagger b$, with the same, now complex, $b_i$, and
\[
E(c)=\int_0^1|f|^2dx-b^\dagger c-c^\dagger b+c^\dagger Ac.
\]
If $Ac_*=b$, then $b^\dagger=c_*^\dagger A$ because $A$ is real and
symmetric, and the same completion of the square gives
\[
E(c_*+d)=\int_0^1|f|^2dx-c_*^\dagger Ac_*+d^\dagger Ad=E(c_*)+d^\dagger Ad,
\]
where $d^\dagger Ad=\int_0^1|\sum_jd_jx^j|^2dx>0$ for $d\ne0$. The same
equations $Ac=b$, now with complex $c$ and $b$, thus give the unique
minimizer.

For $n=5$, Gaussian elimination in exact fractions, without row
interchanges, gives the pivots $1$, $1/12$, $1/180$, $1/2800$, $1/44100$
(Exercise 2.1.5, where the indices start at $1$), and the determinant is
their product:
\[
\det A=\frac1{12\cdot180\cdot2800\cdot44100}=\frac1{266716800000}
=3.7492951325\ldots\times10^{-12}.
\]
\erratum{The Note infers ill-conditioning from the small determinant, but a
small determinant alone is not a condition-number test (rescaling
any well-conditioned matrix can make its determinant small). For this
particular matrix the condition number
$\kappa_2(A)=\sigma_{\max}/\sigma_{\min}\simeq4.7661\times10^5$, the ratio of
the extreme singular values, confirms sensitivity of the equations.}
Because $A$ is symmetric and positive definite, its singular values are its
eigenvalues. Computing them numerically (with any symmetric eigenvalue
routine, e.g., the Jacobi method) gives $\lambda_{\max}\simeq1.56705$ and
$\lambda_{\min}\simeq3.28793\times10^{-6}$, whose ratio is the quoted
$\kappa_2$. To see what this ratio measures, let $b$ change by $\delta b$,
so that $c$ changes by $\delta c=A^{-1}\delta b$. Expanding in orthonormal
eigenvectors of $A$ shows that, for Euclidean lengths,
$|\delta c|\le|\delta b|/\lambda_{\min}$ and
$|b|=|Ac|\le\lambda_{\max}|c|$. Hence
\[
\frac{|\delta c|}{|c|}\le\frac{|\delta b|/\lambda_{\min}}{|b|/\lambda_{\max}}
=\kappa_2(A)\,\frac{|\delta b|}{|b|},
\]
with equality when $c$ is an eigenvector for $\lambda_{\max}$ and
$\delta b$ one for $\lambda_{\min}$. A relative error in $b$ can thus be
magnified by a factor of up to about $5\times10^5$ in the solution $c$.

\Needspace{5\baselineskip}
\par\bigskip
\noindent\textbf{5.1.4}\quad \erratum{Let $F$ and $\varphi_n$ be real, square-integrable functions,
with the $\varphi_n$ orthonormal for a nonnegative weight $w$.}
In place of the expansion of $F(x)$ given by
\[
F(x)=\sum_{n=0}^{\infty}a_{n}\varphi_{n}(x),
\]
with
\[
a_{n}=\int_{a}^{b}F(x)\varphi_{n}(x)w(x)\,dx,
\]
take the finite series approximation
\[
F(x)\approx\sum_{n=0}^{m}c_{n}\varphi_{n}(x).
\]
Show that the mean square error
\[
\int_{a}^{b}\left[F(x)-\sum_{n=0}^{m}c_{n}\varphi_{n}(x)\right]^{2}w(x)\,dx
\]
is minimized by taking $c_{n}=a_{n}$.

\emph{Note.}~The values of the coefficients are independent of the number of
terms in the finite series. This independence is a consequence of
orthogonality and would not hold for a least-squares fit using powers of $x$.
\solution
Use the real orthonormal functions and nonnegative weight implicit in the
stated coefficient formula:
$\int_a^b\varphi_j\varphi_k w\,dx=\delta_{jk}$. Expanding the finite error,
\begin{align*}
E(c)&=\int_a^bF^2w\,dx
 -2\sum_{n=0}^m c_n\int_a^bF\varphi_nw\,dx
 +\sum_{n,k=0}^m c_nc_k\delta_{nk}\\
&=\|F\|^2-2\sum_{n=0}^m c_na_n+\sum_{n=0}^m c_n^2\\
&=\|F\|^2-\sum_{n=0}^m a_n^2+\sum_{n=0}^m(c_n-a_n)^2.
\end{align*}
Here $\|F\|^2=\int_a^bF^2w\,dx$, and the last line uses
$c_n^2-2c_na_n=(c_n-a_n)^2-a_n^2$ to complete the square term by term.
Only the last sum depends on the trial coefficients. It is nonnegative
and vanishes exactly when
\[
\boxed{c_n=a_n\quad(0\le n\le m)},\qquad
E_{\min}=\|F\|^2-\sum_{n=0}^m a_n^2.
\]
Adding another orthonormal function only adds its own square to this
calculation; it does not alter the earlier coefficients. If the full
expansion converges in norm, Parseval's identity
$\|F\|^2=\sum_{n=0}^\infty a_n^2$ gives $E_{\min}=\sum_{n>m}a_n^2$.
For complex functions the corresponding formula uses
$a_n=\int\varphi_n^*Fw\,dx$ and $E=\int|F-\sum c_n\varphi_n|^2w\,dx$;
the last sum is then $\sum|c_n-a_n|^2$, with the same conclusion.

\Needspace{5\baselineskip}
\par\bigskip
\noindent\textbf{5.1.5}\quad From Example 5.1.6,
\[
f(x)=\begin{cases}
\phantom{-}\dfrac{h}{2},&0<x<\pi\\[2ex]
-\dfrac{h}{2},&-\pi<x<0
\end{cases}
=\frac{2h}{\pi}\sum_{n=0}^{\infty}\frac{\sin(2n+1)x}{2n+1}.
\]
\begin{enumerate}
\item[(a)] Show that
  \[
  \int_{-\pi}^{\pi}[f(x)]^{2}\,dx=\frac{\pi}{2}h^{2}
  =\frac{4h^{2}}{\pi}\sum_{n=0}^{\infty}(2n+1)^{-2}.
  \]
  For a finite upper limit this would be Bessel's inequality. For the upper
  limit $\infty$, this is Parseval's identity.
\item[(b)] Verify that
  \[
  \frac{\pi}{2}h^{2}=\frac{4h^{2}}{\pi}\sum_{n=0}^{\infty}(2n+1)^{-2}
  \]
  by evaluating the series.
\end{enumerate}
\emph{Hint.}~The series can be expressed in terms of the Riemann zeta function
$\zeta(2)=\pi^{2}/6$.
\solution
(a) Except at the irrelevant endpoints and jump, $f^2=h^2/4$. Thus
\[
\int_{-\pi}^{\pi} f^2\,dx=2\pi\frac{h^2}{4}=\frac{\pi h^2}{2}.
\]
For positive integers $k,\ell$, the product-to-sum identity gives
\[
\int_{-\pi}^{\pi}\sin kx\sin\ell x\,dx
=\frac12\int_{-\pi}^{\pi}
[\cos(k-\ell)x-\cos(k+\ell)x]dx=\pi\delta_{k\ell}.
\]
The norm squared of the partial sum $f_N$ through $n=N$ is therefore
\begin{align*}
\int_{-\pi}^{\pi}f_N^2\,dx
&=\frac{4h^2}{\pi^2}
 \sum_{n,m=0}^N\frac{\pi\delta_{nm}}{(2n+1)(2m+1)}\\
&=\frac{4h^2}{\pi}\sum_{n=0}^N\frac1{(2n+1)^2}.
\end{align*}
The residual is orthogonal to this partial sum. Indeed, $f\sin(2n+1)x$ is
even, so
\[
\int_{-\pi}^{\pi}f\sin(2n+1)x\,dx
=2\cdot\frac h2\int_0^\pi\sin(2n+1)x\,dx
=h\left[-\frac{\cos(2n+1)x}{2n+1}\right]_0^\pi=\frac{2h}{2n+1},
\]
and hence
\[
\langle f_N|f\rangle=\frac{2h}{\pi}\sum_{n=0}^N\frac1{2n+1}\,\frac{2h}{2n+1}
=\frac{4h^2}{\pi}\sum_{n=0}^N\frac1{(2n+1)^2}=\|f_N\|^2,
\]
i.e., $\langle f_N|f-f_N\rangle=0$. By Pythagoras,
\[
\|f\|^2=\|f_N\|^2+2\langle f_N|f-f_N\rangle+\|f-f_N\|^2
=\|f_N\|^2+\|f-f_N\|^2\ge\|f_N\|^2,
\]
which is Bessel's inequality.
The functions $1,\cos kx,\sin kx$ ($k\ge1$) form a complete orthogonal set
on $[-\pi,\pi]$, so the Fourier series of every square-integrable function
converges to it in the mean-square norm. Here all cosine coefficients
vanish ($f$ is odd), and so do the even-order sine coefficients, since
$\int_0^\pi\sin2mx\,dx=0$; the partial sums of the Fourier series of $f$
are therefore the $f_N$, and $\|f-f_N\|\to0$. Letting $N\to\infty$ in the
Pythagoras identity proves the stated Parseval equality.

(b) Split the positive integers into odd and even ones:
\[
\sum_{n=0}^\infty\frac1{(2n+1)^2}
=\sum_{k=1}^\infty\frac1{k^2}
 -\sum_{k=1}^\infty\frac1{(2k)^2}
=\left(1-\frac14\right)\zeta(2)=\frac{\pi^2}{8}.
\]
Substitution yields
$(4h^2/\pi)(\pi^2/8)=\pi h^2/2$, an independent check of the equality.

\Needspace{5\baselineskip}
\par\bigskip
\noindent\textbf{5.1.6}\quad Derive the Schwarz inequality \erratum{for real square-integrable functions} from the identity
\[
\left[\int_{a}^{b}f(x)g(x)\,dx\right]^{2}
=\int_{a}^{b}f^{2}(x)\,dx\int_{a}^{b}g^{2}(x)\,dx
-\frac{1}{2}\int_{a}^{b}dx\int_{a}^{b}dy\,
[f(x)g(y)-f(y)g(x)]^{2}.
\]
\solution
\erratum{The qualification is needed for the derivation: square
integrability makes every integral below finite, and for real functions
the integrand of the double integral is a nonnegative square. (The
identity is pure algebra and also holds for complex functions, but then
that integrand need not be nonnegative.)}
To check the identity, expand its double integral. Fubini's theorem
applies since $\int_a^bf^2dx$ and $\int_a^bg^2dx$ are finite and
$\int_a^b|fg|\,dx<\infty$. The latter fact follows from
$2|fg|\le f^2+g^2$, so it does not assume the Schwarz inequality we
are proving:
\begin{align*}
D&=\int_a^b\!\int_a^b[f(x)g(y)-f(y)g(x)]^2\,dx\,dy\\
&=\int_a^b\!\int_a^b
 \{f(x)^2g(y)^2+f(y)^2g(x)^2
          -2f(x)g(x)f(y)g(y)\}\,dx\,dy\\
&=2\left(\int_a^bf^2dx\right)\left(\int_a^bg^2dx\right)
 -2\left(\int_a^bfg\,dx\right)^2.
\end{align*}
In the last step each term, a function of $x$ times a function of $y$,
integrates to a product of single integrals, e.g.,
$\int\!\!\int f(x)^2g(y)^2\,dx\,dy=\int f^2dx\int g^2dy$; the second term
gives the same product with $x$ and $y$ renamed, and the third gives
$-2(\int fg\,dx)(\int fg\,dy)$. Dividing by $2$ gives
\[
\frac12D=\int_a^bf^2dx\int_a^bg^2dx-\left(\int_a^bfg\,dx\right)^2,
\]
which is the identity.
Since the original integrand in $D$ is a square, $D\ge0$, and the
identity gives $(\int fg\,dx)^2\le\int f^2dx\int g^2dx$. Taking
nonnegative square roots proves
\[
\boxed{\left|\int_a^bfg\,dx\right|
\le\left(\int_a^bf^2dx\right)^{1/2}
   \left(\int_a^bg^2dx\right)^{1/2}.}
\]
If $\int g^2dx>0$, set $\lambda=(\int fg)/(\int g^2)$. Then
\begin{align*}
\int_a^b(f-\lambda g)^2dx
&=\int f^2dx-2\lambda\int fg\,dx+\lambda^2\int g^2dx\\
&=\int f^2dx-\frac{(\int fg\,dx)^2}{\int g^2dx}.
\end{align*}
The right side vanishes exactly in the case of equality, and the left
side exactly when $f=\lambda g$ almost everywhere. Thus equality holds
exactly when $f=\lambda g$ almost everywhere; the zero-function cases also
give equality. For complex functions the
Schwarz inequality is $|\int f^*g|^2\le\int|f|^2\int|g|^2$;
ordinary squares in the displayed real identity must be replaced by
appropriate modulus squares.

\Needspace{5\baselineskip}
\par\bigskip
\noindent\textbf{5.1.7}\quad Starting from
\[
I=\left\langle f-\sum_{i}a_{i}\varphi_{i}\,\middle|\,
f-\sum_{j}a_{j}\varphi_{j}\right\rangle\ge0,
\]
derive Bessel's inequality,
$\langle f|f\rangle\ge\sum_{n}|a_{n}|^{2}$.
\erratum{Here the $\varphi_n$ are orthonormal and
$a_n=\langle\varphi_n|f\rangle$; begin with finite partial sums.}
\solution
\erratum{The added hypotheses are needed: for $f=\varphi_0$, the
arbitrary choice $a_0=2$ (all other $a_n=0$) gives
$\sum|a_n|^2=4>\langle f|f\rangle=1$, and $a_n=\langle\varphi_n|f\rangle$
with the nonorthogonal choice $\varphi_1=\varphi_0$ gives
$|a_0|^2+|a_1|^2=2>1$. Finite sums come first because the convergence of
$\sum_ia_i\varphi_i$, which the printed $I$ presupposes, is not known in
advance.}
Take a finite set of orthonormal functions first, and use the projection
coefficients $a_j=\langle\varphi_j|f\rangle$. The scalar product is
conjugate-linear in its first argument, so
\begin{align*}
I_N&=\left\langle f-\sum_{i=0}^Na_i\varphi_i\,\middle|\,
 f-\sum_{j=0}^Na_j\varphi_j\right\rangle\\
&=\langle f|f\rangle
 -\sum_j a_j\langle f|\varphi_j\rangle
 -\sum_i a_i^*\langle\varphi_i|f\rangle
 +\sum_{i,j}a_i^*a_j\delta_{ij}\\
&=\|f\|^2-\sum_j|a_j|^2-\sum_i|a_i|^2
 +\sum_i|a_i|^2
=\|f\|^2-\sum_{i=0}^N|a_i|^2\ge0.
\end{align*}
The passage to the last line uses $\langle\varphi_i|f\rangle=a_i$ and
$\langle f|\varphi_j\rangle=\langle\varphi_j|f\rangle^*=a_j^*$.
The partial sums of the nonnegative series $\sum|a_i|^2$ increase and
are bounded by $\|f\|^2$. Taking their limit gives
\[
\boxed{\sum_{i=0}^\infty|a_i|^2\le\langle f|f\rangle.}
\]
No completeness assumption on the orthonormal set is needed for this
inequality. To check the equality case, let $S_N=\sum_{i=0}^Na_i\varphi_i$.
For $M>N$, orthonormality gives
$\|S_M-S_N\|^2=\sum_{i=N+1}^M|a_i|^2$, which tends to $0$ as $N\to\infty$
(a tail of the convergent series), so the $S_N$ form a Cauchy sequence.
A Hilbert space is complete in the sense that every Cauchy sequence
converges in it (a property of the space, not of the orthonormal set), so
$S_N\to S$ for some $S$. This $S$ lies in the closed span of the
$\varphi_i$, the set of all norm limits of finite linear combinations of
the $\varphi_i$.
Letting $N\to\infty$ in the residual identity
$\|f-S_N\|^2=I_N=\|f\|^2-\sum_{i=0}^N|a_i|^2$ (the norm is continuous) gives
\[
\|f-S\|^2=\|f\|^2-\sum_{i=0}^\infty|a_i|^2.
\]
Thus equality implies $f=S$ in that closed span. Conversely, if $f$ is in
the closed span, the residual $f-S$ belongs to it and is orthogonal to
every $\varphi_k$, since
$\langle\varphi_k|f-S\rangle=a_k-\lim_N\langle\varphi_k|S_N\rangle=a_k-a_k=0$;
by linearity and continuity it is orthogonal to the entire closed span.
Its squared norm is therefore zero, proving equality.

\Needspace{5\baselineskip}
\par\bigskip
\noindent\textbf{5.1.8}\quad Expand the function $\sin\pi x$ in a series of
functions $\varphi_{i}$ that are orthogonal (but not normalized) on the range
$0\le x\le1$ when the scalar product has definition
\[
\langle f|g\rangle=\int_{0}^{1}f^{*}(x)g(x)\,dx.
\]
Keep the first four terms of the expansion. The first four $\varphi_{i}$ are:
\[
\varphi_{0}=1,\quad\varphi_{1}=2x-1,\quad\varphi_{2}=6x^{2}-6x+1,\quad
\varphi_{3}=20x^{3}-30x^{2}+12x-1.
\]
\emph{Note.}~The integrals that are needed are the subject of Example 1.10.5.
\solution
For an orthogonal but unnormalized family,
$c_n=\langle\varphi_n|f\rangle/\langle\varphi_n|\varphi_n\rangle$.
Direct polynomial integration gives the four squared norms
\[
\int_0^1\varphi_n^2dx=1,\ \frac13,\ \frac15,\ \frac17
\qquad(n=0,1,2,3).
\]
For the nonconstant squared norms, square the polynomials,
\begin{align*}
\varphi_1^2&=4x^2-4x+1,\qquad
\varphi_2^2=36x^4-72x^3+48x^2-12x+1,\\
\varphi_3^2&=400x^6-1200x^5+1380x^4-760x^3+204x^2-24x+1,
\end{align*}
and integrate term by term with $\int_0^1x^k\,dx=1/(k+1)$:
\begin{align*}
\|\varphi_1\|^2&=\frac43-\frac42+1=\frac13,\\
\|\varphi_2\|^2&=\frac{36}5-\frac{72}4+\frac{48}3-\frac{12}2+1=\frac15,\\
\|\varphi_3\|^2&=\frac{400}7-\frac{1200}6+\frac{1380}5
 -\frac{760}4+\frac{204}3-\frac{24}2+1=\frac17.
\end{align*}
To compute all numerators, put $I_j=\int_0^1x^j\sin\pi x\,dx$.
The first two integrals are
\[
I_0=\frac2\pi,\qquad
I_1=\left[-\frac{x\cos\pi x}{\pi}
 +\frac{\sin\pi x}{\pi^2}\right]_0^1=\frac1\pi.
\]
Two integrations by parts, for $j\ge2$, give
\begin{align*}
I_j&=\left[-\frac{x^j\cos\pi x}{\pi}\right]_0^1
 +\frac{j}{\pi}\int_0^1x^{j-1}\cos\pi x\,dx\\
&=\frac1\pi+\frac{j}{\pi}\left(
 \left[\frac{x^{j-1}\sin\pi x}{\pi}\right]_0^1
 -\frac{j-1}{\pi}\int_0^1x^{j-2}\sin\pi x\,dx\right)\\
&=\frac1\pi-\frac{j(j-1)}{\pi^2}I_{j-2},
\end{align*}
since $\cos\pi=-1$ and $\sin\pi=\sin0=0$. Hence
\[
I_2=\frac1\pi-\frac2{\pi^2}I_0=\frac1\pi-\frac4{\pi^3},\qquad
I_3=\frac1\pi-\frac6{\pi^2}I_1=\frac1\pi-\frac6{\pi^3}.
\]
Substituting,
\begin{align*}
\langle\varphi_0|f\rangle&=I_0=2/\pi,\\
\langle\varphi_1|f\rangle&=2I_1-I_0=0,\\
\langle\varphi_2|f\rangle&=6I_2-6I_1+I_0
=2/\pi-24/\pi^3,\\
\langle\varphi_3|f\rangle&=20I_3-30I_2+12I_1-I_0
=\frac{20-30+12-2}{\pi}+\frac{-120+120}{\pi^3}=0.
\end{align*}
Dividing by the squared norms $\|\varphi_0\|^2=1$ and
$\|\varphi_2\|^2=1/5$ gives
\[
c_0=\frac2\pi,\qquad
c_2=5\Bigl(\frac2\pi-\frac{24}{\pi^3}\Bigr)=\frac{10(\pi^2-12)}{\pi^3}
\approx-0.6871,
\]
while $c_1=c_3=0$. The requested four-term projection is therefore
\[
\boxed{\sin\pi x\ \approx\ \frac2\pi\varphi_0
 +0\varphi_1+\frac{10(\pi^2-12)}{\pi^3}\varphi_2+0\varphi_3.}
\]
The two zero coefficients can also be checked by reflection:
\[
\sin\pi(1-x)=\sin\pi x,\qquad\text{whereas}\qquad
\varphi_1(1-x)=-\varphi_1(x),\quad \varphi_3(1-x)=-\varphi_3(x).
\]
The substitution $x\to1-x$, which maps
$[0,1]$ onto itself and so leaves the value of an integral over $[0,1]$
unchanged, therefore turns each of $\int_0^1\varphi_1\sin\pi x\,dx$ and
$\int_0^1\varphi_3\sin\pi x\,dx$ into its own negative; both integrals
vanish.
Because each $c_n$ is a projection coefficient, the residual
$\sin\pi x-\sum_{n=0}^3c_n\varphi_n$ has zero scalar product with each of
the four retained functions.

\Needspace{5\baselineskip}
\par\bigskip
\noindent\textbf{5.1.9}\quad Expand the function $e^{-x}$ in Laguerre
polynomials $L_{n}(x)$, which are orthonormal on the range $0\le x<\infty$
with scalar product
\[
\langle f|g\rangle=\int_{0}^{\infty}f^{*}(x)g(x)e^{-x}\,dx.
\]
Keep the first four terms of the expansion. The first four $L_{n}(x)$ are
\[
L_{0}=1,\quad L_{1}=1-x,\quad L_{2}=\frac{2-4x+x^{2}}{2},\quad
L_{3}=\frac{6-18x+9x^{2}-x^{3}}{6}.
\]
\solution
The weight contributes an additional factor $e^{-x}$, so
\[
a_n=\langle L_n|e^{-x}\rangle=\int_0^\infty L_n(x)e^{-2x}dx.
\]
Set $J_k=\int_0^\infty x^ke^{-2x}dx$. Then $J_0=1/2$, and for $k\ge1$
integration by parts gives
\[
J_k=\left[-\frac{x^ke^{-2x}}{2}\right]_0^\infty
 +\frac k2\int_0^\infty x^{k-1}e^{-2x}dx=\frac k2J_{k-1},
\]
hence $J_k=k!/2^{k+1}$, i.e., $(J_0,J_1,J_2,J_3)=(1/2,1/4,1/4,3/8)$.
Inserting the powers of $x$ in each $L_n$, the four coefficients are
\begin{align*}
a_0&=J_0=\frac12,\\
a_1&=J_0-J_1=\frac12-\frac14=\frac14,\\
a_2&=J_0-2J_1+\frac12J_2
=\frac12-\frac12+\frac18=\frac18,\\
a_3&=J_0-3J_1+\frac32J_2-\frac16J_3
=\frac12-\frac34+\frac38-\frac1{16}=\frac1{16}.
\end{align*}
Consequently
\[
\boxed{e^{-x}\approx\frac12L_0+\frac14L_1+\frac18L_2+\frac1{16}L_3
=\frac{15}{16}-\frac{11x}{16}+\frac{5x^2}{32}-\frac{x^3}{96}.}
\]
Collecting powers of $x$: the constant is
$\frac12+\frac14+\frac18+\frac1{16}=\frac{15}{16}$, the coefficient of $x$ is
$-\frac14-\frac14-\frac3{16}=-\frac{11}{16}$, that of $x^2$ is
$\frac1{16}+\frac3{32}=\frac5{32}$, and that of $x^3$ is $-\frac1{96}$.
As a check of both the weight and the normalization, the residual identity
of Exercise 5.1.7, with $\|e^{-x}\|^2=\int_0^\infty e^{-2x}e^{-x}dx$, gives
the squared error norm of the four-term sum $f_3$:
\begin{align*}
\|e^{-x}-f_3\|^2
&=\int_0^\infty e^{-3x}dx-\sum_{n=0}^3|a_n|^2\\
&=\frac13-\left(\frac14+\frac1{16}+\frac1{64}+\frac1{256}\right)
=\frac13-\frac{85}{256}=\frac1{768}.
\end{align*}
Expanding the square gives the same value:
\begin{align*}
\int_0^\infty(e^{-x}-f_3)^2e^{-x}dx
&=\int_0^\infty e^{-3x}dx-2\int_0^\infty f_3e^{-2x}dx
 +\int_0^\infty f_3^2e^{-x}dx\\
&=\frac13-2\cdot\frac{85}{256}+\frac{85}{256}=\frac1{768},
\end{align*}
where the polynomial form of $f_3$ and the $J_k$ give
$\int_0^\infty f_3e^{-2x}dx
=\frac{15}{32}-\frac{11}{64}+\frac5{128}-\frac1{256}=\frac{85}{256}$, and
orthonormality gives
$\int_0^\infty f_3^2e^{-x}dx=\sum_{n=0}^3a_n^2=\frac{85}{256}$.

\Needspace{5\baselineskip}
\par\bigskip
\noindent\textbf{5.1.10}\quad The explicit form of a function $f$ is not
known, but the coefficients $a_{n}$ of its expansion in the orthonormal set
$\varphi_{n}$ are available. Assuming that the $\varphi_{n}$ and the members
of another orthonormal set, $\chi_{n}$, are available, use Dirac notation to
obtain a formula for the coefficients for the expansion of $f$ in the
$\chi_{n}$ set.
\solution
Write $|f\rangle=\sum_n a_n|\varphi_n\rangle$. Taking the scalar product
with the $m$th new basis vector gives
\[
\boxed{b_m=\langle\chi_m|f\rangle
=\sum_n\langle\chi_m|\varphi_n\rangle a_n.}
\]
Thus only overlaps between known basis functions are required. In matrix
notation $b=Ua$, where $U_{mn}=\langle\chi_m|\varphi_n\rangle$.
For an infinite expansion the interchange of $\langle\chi_m|$ and the sum
follows by taking scalar products with its norm-convergent partial sums,
as in Exercise 5.1.1.

If both sets are complete in the same space, then
$(U^\dagger)_{jm}=U_{mj}^*=\langle\varphi_j|\chi_m\rangle$, and inserting
the resolution of the identity $\sum_m|\chi_m\rangle\langle\chi_m|=1$ gives
\[
(U^\dagger U)_{jk}
=\sum_m\langle\varphi_j|\chi_m\rangle
           \langle\chi_m|\varphi_k\rangle
=\langle\varphi_j|\varphi_k\rangle=\delta_{jk}.
\]
Similarly, $\sum_k|\varphi_k\rangle\langle\varphi_k|=1$ gives
\[
(UU^\dagger)_{mn}
=\sum_k\langle\chi_m|\varphi_k\rangle
           \langle\varphi_k|\chi_n\rangle
=\langle\chi_m|\chi_n\rangle=\delta_{mn},
\]
i.e., $UU^\dagger=I$. If the $\chi_m$ do not span $f$, the same
formula instead gives the coefficients of its orthogonal projection;
it cannot reconstruct a component outside their closed span.

\Needspace{5\baselineskip}
\par\bigskip
\noindent\textbf{5.1.11}\quad Using conventional vector notation, evaluate
$\sum_{j}|\hat{\mathbf{e}}_{j}\rangle\langle\hat{\mathbf{e}}_{j}|
\mathbf{a}\rangle$, where $\mathbf{a}$ is an arbitrary vector in the space
spanned by the \erratum{orthonormal} $\hat{\mathbf{e}}_{j}$.
\solution
For the usual orthonormal unit vectors,
$\langle\hat{\mathbf e}_j|\mathbf a\rangle
=\hat{\mathbf e}_j\cdot\mathbf a=a_j$. Therefore
\[
\boxed{\sum_j|\hat{\mathbf e}_j\rangle
 \langle\hat{\mathbf e}_j|\mathbf a\rangle
=\sum_j\hat{\mathbf e}_j(\hat{\mathbf e}_j\cdot\mathbf a)
=\sum_j a_j\hat{\mathbf e}_j=\mathbf a.}
\]
In ordinary dyadic notation this is
$(\sum_j\hat{\mathbf e}_j\hat{\mathbf e}_j^{\mathsf T})\mathbf a=I\mathbf a$.
\erratum{Orthonormality is needed here; unit length and spanning alone are
insufficient.} For example, the unit vectors $(1,0)$ and
$(1,1)/\sqrt2$ span $\mathbb R^2$, but the displayed sum applied to
$(1,0)$ is
\[
(1,0)\,[(1,0)\cdot(1,0)]
+\frac{(1,1)}{\sqrt2}\Bigl[\frac{(1,1)}{\sqrt2}\cdot(1,0)\Bigr]
=(1,0)+\frac12(1,1)=\Bigl(\frac32,\frac12\Bigr),
\]
not $(1,0)$.

\Needspace{5\baselineskip}
\par\bigskip
\noindent\textbf{5.1.12}\quad Letting
$\mathbf{a}=a_{1}\hat{\mathbf{e}}_{1}+a_{2}\hat{\mathbf{e}}_{2}$ and
$\mathbf{b}=b_{1}\hat{\mathbf{e}}_{1}+b_{2}\hat{\mathbf{e}}_{2}$ be vectors in
$\mathbb{R}^{2}$, for what values of $k$, if any, is
\[
\langle\mathbf{a}|\mathbf{b}\rangle
=a_{1}b_{1}-a_{1}b_{2}-a_{2}b_{1}+ka_{2}b_{2}
\]
a valid definition of a scalar product?
\solution
The proposed form is bilinear and symmetric for every real $k$, because
it equals $(a_1,a_2)\,G\,(b_1,b_2)^T$ with the matrix
\[
G=\begin{pmatrix}1&-1\\-1&k\end{pmatrix}=G^T.
\]
It remains to test positive definiteness. Completing the square,
\[
\langle\mathbf a|\mathbf a\rangle
=a_1^2-2a_1a_2+ka_2^2=(a_1-a_2)^2+(k-1)a_2^2.
\]
If $k>1$, both terms are nonnegative and their sum can vanish only if
$a_2=0$ and then $a_1=0$. If $k=1$, the nonzero vector $(1,1)$ has
zero norm. If $k<1$, this same vector has squared norm $k-1<0$.
Thus all scalar-product axioms hold exactly for
\[
\boxed{k>1.}
\]
The determinant check $\det G=k-1>0$, together with $G_{11}=1>0$
(Sylvester's criterion for positive definiteness), agrees with the direct
square-completion argument.

\Needspace{9\baselineskip}
\subsection*{5.2 Gram-Schmidt Orthogonalization}

For the Gram-Schmidt constructions in Exercises 5.2.1 through 5.2.6, use a
scalar product of the form given in Eq.~(5.7) with the specified interval and
weight.

\Needspace{5\baselineskip}
\par\bigskip
\noindent\textbf{5.2.1}\quad Following the Gram-Schmidt procedure, construct a
set of polynomials $P_{n}^{*}(x)$ orthogonal (unit weighting factor) over the
range $[0,1]$ from the set $[1,x,x^{2},\ldots]$. Scale so that
$P_{n}^{*}(1)=1$.
\par\nobreak\hfill\emph{ANS.}\enspace
$\begin{aligned}
\erratum{P_{0}^{*}(x)}&=1,\\
P_{1}^{*}(x)&=2x-1,\\
P_{2}^{*}(x)&=6x^{2}-6x+1,\\
P_{3}^{*}(x)&=20x^{3}-30x^{2}+12x-1.
\end{aligned}$\par
These are the first four shifted Legendre polynomials.

\emph{Note.}~The ``*'' is the standard notation for ``shifted'': $[0,1]$
instead of $[-1,1]$. It does not mean complex conjugate.
\solution
\erratum{The first answer on textbook p.~273 is printed $P_n^*(x)=1$;
its subscript must be $0$, because $P_n^*$ has degree $n$, so the
constant answer is $P_0^*$.}
It is convenient first to construct monic orthogonal polynomials $q_n$:
\[
q_n=x^n-\sum_{j=0}^{n-1}
 \frac{\int_0^1q_jx^n\,dx}{\int_0^1q_j^2\,dx}\,q_j,
\qquad P_n^*=\frac{q_n}{q_n(1)}.
\]
All the required integrals reduce to $\int_0^1x^kdx=1/(k+1)$.
Start with $q_0=1$, $\|q_0\|^2=1$. Then
\[
q_1=x-\frac12,\qquad
\|q_1\|^2=\frac13-\frac12+\frac14=\frac1{12}.
\]
For the next polynomial,
\[
\langle q_0|x^2\rangle=\frac13,\qquad
\langle q_1|x^2\rangle=\frac14-\frac16=\frac1{12}.
\]
Dividing by $\|q_0\|^2=1$ and $\|q_1\|^2=1/12$, the coefficients are $1/3$
and $1$, so
\[
q_2=x^2-\frac13-q_1=x^2-x+\frac16,
\]
\[
q_2^2=x^4-2x^3+\frac43x^2-\frac13x+\frac1{36},\qquad
\|q_2\|^2
=\frac15-\frac12+\frac49-\frac16+\frac1{36}=\frac1{180}.
\]
For $x^3$ the three projection numerators are
\[
\langle q_0|x^3\rangle=\frac14,\qquad
\langle q_1|x^3\rangle=\frac15-\frac18=\frac3{40},\qquad
\langle q_2|x^3\rangle=\frac16-\frac15+\frac1{24}=\frac1{120}.
\]
Dividing by $1,1/12,1/180$, respectively, gives the coefficients
$1/4$, $9/10$, $3/2$ and
\begin{align*}
q_3&=x^3-\frac14-\frac9{10}\left(x-\frac12\right)
 -\frac32\left(x^2-x+\frac16\right)\\
&=x^3-\frac32x^2+\frac35x-\frac1{20},
\end{align*}
since the coefficient of $x$ is $-\frac9{10}+\frac32=\frac35$ and the
constant is $-\frac14+\frac9{20}-\frac14=-\frac1{20}$.
The endpoint values are $q_0(1)=1$, $q_1(1)=1/2$,
$q_2(1)=1/6$, and $q_3(1)=(20-30+12-1)/20=1/20$. Hence
\[
\boxed{P_0^*=1,\quad P_1^*=2x-1,\quad
P_2^*=6x^2-6x+1,\quad P_3^*=20x^3-30x^2+12x-1.}
\]
The construction extends to the complete set. A formula with the same
orthogonality and endpoint normalization is
\[
\boxed{P_n^*(x)=\frac1{n!}\frac{d^n}{dx^n}[x^n(x-1)^n]
=\sum_{k=0}^n(-1)^{n-k}\binom nk\binom{n+k}k x^k.}
\]
The finite sum follows by first expanding and then differentiating:
\[
x^n(x-1)^n=\sum_{k=0}^n(-1)^{n-k}\binom nk x^{n+k},\qquad
\frac1{n!}\frac{d^n x^{n+k}}{dx^n}
=\frac{(n+k)!}{n!k!}x^k=\binom{n+k}k x^k.
\]
To verify orthogonality, multiply the derivative formula by $x^j$, $j<n$, and
integrate by parts $n$ times. All boundary terms vanish because
$x^n(x-1)^n$ and its first $n-1$ derivatives vanish at both endpoints;
the remaining integral vanishes because $d^nx^j/dx^n=0$.
At $x=1$, use the Leibniz rule
\[
\frac{d^n}{dx^n}[x^n(x-1)^n]=\sum_{k=0}^n\binom nk
\frac{d^kx^n}{dx^k}\,\frac{d^{n-k}(x-1)^n}{dx^{n-k}}.
\]
For $k\ge1$ the last factor is $(x-1)^k\,n!/k!$, which vanishes at
$x=1$, so only the term differentiating $(x-1)^n$ all $n$ times ($k=0$)
survives; it equals $1^n\,n!=n!$, so $P_n^*(1)=1$.
Finally, these are the Gram--Schmidt polynomials: the polynomial given by
the derivative formula and the Gram--Schmidt $P_n^*$ (orthogonal to
$P_0^*,\ldots,P_{n-1}^*$, which span the polynomials of degree below $n$)
both have degree $n$, are orthogonal to every polynomial of lower degree,
and equal $1$ at $x=1$. If $R$ and $S$ are two such polynomials,
$R-\lambda S$, with $\lambda$ the ratio of their leading coefficients, has
lower degree and is orthogonal to itself, hence zero; $x=1$ then gives
$\lambda=1$, so $R=S$.

A useful check is the norm formula. The leading coefficient is
$\binom{2n}n$ and lower powers are orthogonal to $P_n^*$, so
\begin{align*}
\int_0^1(P_n^*)^2dx
&=\binom{2n}n\int_0^1x^nP_n^*(x)dx\\
&=\binom{2n}n\int_0^1x^n(1-x)^n dx
=\frac{(2n)!}{(n!)^2}\frac{(n!)^2}{(2n+1)!}
=\frac1{2n+1}.
\end{align*}
The middle equality again uses $n$ integrations by parts with vanishing
boundary terms, now moving all derivatives onto $x^n$:
\[
\int_0^1x^nP_n^*dx
=\frac{(-1)^n}{n!}\int_0^1\frac{d^nx^n}{dx^n}\,x^n(x-1)^n\,dx
=(-1)^n\int_0^1x^n(x-1)^ndx=\int_0^1x^n(1-x)^ndx.
\]
To evaluate the
last integral without a table, set $J_{a,b}=\int_0^1x^a(1-x)^b dx$.
For integers $a\ge0$, $b\ge1$, integration by parts gives
\[
J_{a,b}=\left[\frac{x^{a+1}(1-x)^b}{a+1}\right]_0^1
 +\frac b{a+1}J_{a+1,b-1}=\frac b{a+1}J_{a+1,b-1}.
\]
Repeating $b$ times and using $J_{a+b,0}=1/(a+b+1)$ yields
\[
J_{a,b}=\frac b{a+1}\,\frac{b-1}{a+2}\cdots\frac1{a+b}\,\frac1{a+b+1}
=\frac{a!\,b!}{(a+b+1)!};
\]
take $a=b=n$.

\Needspace{5\baselineskip}
\par\bigskip
\noindent\textbf{5.2.2}\quad Apply the Gram-Schmidt procedure to form the
first three Laguerre polynomials:
\[
u_{n}(x)=x^{n},\quad n=0,1,2,\ldots,\quad0\le x<\infty,\quad w(x)=e^{-x}.
\]
The conventional normalization is
\[
\int_{0}^{\infty}L_{m}(x)L_{n}(x)e^{-x}\,dx=\delta_{mn}.
\]
\par\nobreak\hfill\emph{ANS.}\enspace
$L_{0}=1$,\quad$L_{1}=(1-x)$,\quad$L_{2}=\dfrac{2-4x+x^{2}}{2}$.
\solution
All scalar products reduce to the moments
$M_k=\int_0^\infty x^ke^{-x}dx$. For $k\ge1$, integration by parts gives
\[
M_k=[-x^ke^{-x}]_0^\infty+kM_{k-1}=kM_{k-1},\qquad M_0=1.
\]
The boundary term is zero, so repeated application gives
\[
M_k=k!,\qquad
M_0=1,\ M_1=1,\ M_2=2,\ M_3=6,\ M_4=24.
\]
Begin with $q_0=1$ and subtract its projection from $x$:
\[
q_1=x-\frac{M_1}{M_0}=x-1,\qquad
\|q_1\|^2=M_2-2M_1+M_0=1.
\]
For $x^2$, the projections onto $q_0$ and $q_1$ have coefficients
\[
\frac{\langle q_0|x^2\rangle}{\|q_0\|^2}=\frac{M_2}{M_0}=2,
\qquad
\frac{\langle q_1|x^2\rangle}{\|q_1\|^2}=M_3-M_2=4.
\]
Therefore
\[
q_2=x^2-2-4(x-1)=x^2-4x+2.
\]
Since $q_2^2=x^4-8x^3+20x^2-16x+4$,
\[
\|q_2\|^2=M_4-8M_3+20M_2-16M_1+4M_0
=24-48+40-16+4=4.
\]
Choose the conventional signs $L_n(0)>0$. Dividing by the norms gives
\[
\boxed{L_0=1,\qquad L_1=-q_1=1-x,\qquad
L_2=q_2/2=1-2x+\frac{x^2}{2}.}
\]
Gram--Schmidt makes the cross integrals zero, and the displayed norm
calculations make all three diagonal integrals one. Normalization alone
leaves each polynomial's overall sign free.

\Needspace{5\baselineskip}
\par\bigskip
\noindent\textbf{5.2.3}\quad You are given
\begin{enumerate}
\item[(a)] a set of functions $u_{n}(x)=x^{n}$, $n=0,1,2,\ldots$,
\item[(b)] an interval $(0,\infty)$,
\item[(c)] a weighting function $w(x)=xe^{-x}$.
\end{enumerate}
Use the Gram-Schmidt procedure to construct the first three orthonormal
functions from the set $u_{n}(x)$ for this interval and this weighting
function.
\par\nobreak\hfill\emph{ANS.}\enspace
$\varphi_{0}(x)=1$,\quad$\varphi_{1}(x)=(x-2)/\sqrt{2}$,\quad
$\varphi_{2}(x)=\erratum{\dfrac{x^{2}-6x+6}{2\sqrt{3}}}$.
\solution
With the extra factor $x$ in the weight, the moments are those of
Exercise 5.2.2 with $k$ raised by one:
\[
M_k=\int_0^\infty x^kxe^{-x}dx=(k+1)!,
\quad (M_0,M_1,M_2,M_3,M_4)=(1,2,6,24,120).
\]
The first function is $q_0=\varphi_0=1$. For the second,
\[
q_1=x-\frac{M_1}{M_0}=x-2,\qquad
\|q_1\|^2=6-4(2)+4(1)=2,
\quad \varphi_1=\frac{x-2}{\sqrt2}.
\]
For the third polynomial,
\[
\langle q_0|x^2\rangle=M_2=6,\qquad
\frac{\langle q_1|x^2\rangle}{\|q_1\|^2}
=\frac{M_3-2M_2}{2}=\frac{24-12}{2}=6.
\]
Consequently
\[
q_2=x^2-6-6(x-2)=x^2-6x+6.
\]
Its square is $x^4-12x^3+48x^2-72x+36$, so
\[
\|q_2\|^2=120-12(24)+48(6)-72(2)+36=12.
\]
Thus
\[
\boxed{\varphi_0=1,\qquad\varphi_1=\frac{x-2}{\sqrt2},\qquad
\varphi_2=\frac{x^2-6x+6}{2\sqrt3}.}
\]
\erratum{The printed slash notation for $\varphi_2$ means division by the
entire quantity $2\sqrt3$; writing a fraction removes the ambiguity.}
As a direct orthogonality check,
$\langle1|q_2\rangle=M_2-6M_1+6M_0=6-12+6=0$ and, since
$(x-2)(x^2-6x+6)=x^3-8x^2+18x-12$,
$\langle q_1|q_2\rangle=24-8(6)+18(2)-12=0$.

\Needspace{5\baselineskip}
\par\bigskip
\noindent\textbf{5.2.4}\quad Using the Gram-Schmidt orthogonalization
procedure, construct the lowest three Hermite polynomials:
\[
u_{n}(x)=x^{n},\quad n=0,1,2,\ldots,\quad-\infty<x<\infty,\quad
w(x)=e^{-x^{2}}.
\]
For this set of polynomials the usual normalization is
\[
\int_{-\infty}^{\infty}H_{m}(x)H_{n}(x)w(x)\,dx
=\delta_{mn}2^{m}m!\,\pi^{1/2}.
\]
\par\nobreak\hfill\emph{ANS.}\enspace
$H_{0}=1$,\quad$H_{1}=2x$,\quad$H_{2}=4x^{2}-2$.
\solution
Odd moments vanish by parity. For $j\ge1$, use
$x e^{-x^2}dx=-\tfrac12d(e^{-x^2})$ in integration by parts:
\[
M_{2j}=-\frac12[x^{2j-1}e^{-x^2}]_{-\infty}^{\infty}
 +\frac{2j-1}{2}\int_{-\infty}^{\infty}x^{2j-2}e^{-x^2}dx.
\]
The boundary term vanishes. The even-moment recurrence is therefore
\[
M_{2j}=\int_{-\infty}^{\infty}x^{2j}e^{-x^2}dx
=\frac{2j-1}{2}M_{2j-2},\qquad M_0=\sqrt\pi.
\]
Thus $M_2=\sqrt\pi/2$ and $M_4=3\sqrt\pi/4$.
Start with $q_0=1$. The projection of $x$ onto $1$ vanishes, so $q_1=x$.
For $x^2$ the projection onto $x$ vanishes, and
\[
q_2=x^2-\frac{M_2}{M_0}=x^2-\frac12.
\]
Their squared norms, using $q_2^2=x^4-x^2+\frac14$, are
\[
\|q_0\|^2=\sqrt\pi,\qquad
\|q_1\|^2=M_2=\frac{\sqrt\pi}{2},\qquad
\|q_2\|^2=M_4-M_2+\frac14M_0
=\Bigl(\frac34-\frac12+\frac14\Bigr)\sqrt\pi=\frac{\sqrt\pi}{2}.
\]
The required squared norms $2^nn!\,\pi^{1/2}$ for degrees $0,1,2$ are
respectively $\sqrt\pi,2\sqrt\pi,8\sqrt\pi$.
The needed positive scale factors $s_n$, with $s_n^2\|q_n\|^2$ equal to
these, are therefore
$1$, $\sqrt{2/(1/2)}=2$, and $\sqrt{8/(1/2)}=4$. Hence
\[
\boxed{H_0=1,\qquad H_1=2x,\qquad H_2=4x^2-2.}
\]
Parity verifies the odd cross integrals, while
$\langle H_0|H_2\rangle=4M_2-2M_0=0$ checks the remaining pair.

\Needspace{5\baselineskip}
\par\bigskip
\noindent\textbf{5.2.5}\quad Use the Gram-Schmidt orthogonalization scheme to
construct the first three Chebyshev polynomials (type I):
\[
u_{n}(x)=x^{n},\quad n=0,1,2,\ldots,\quad-1\le x\le1,\quad
w(x)=(1-x^{2})^{-1/2}.
\]
Take the normalization
\[
\int_{-1}^{1}T_{m}(x)T_{n}(x)w(x)\,dx
=\delta_{mn}\begin{cases}\pi,&m=n=0,\\[1ex]
\dfrac{\pi}{2},&m=n\ge1.
\end{cases}
\]
\emph{Hint.}~The needed integrals are given in Exercise 13.3.2.
\par\nobreak\hfill\emph{ANS.}\enspace
$T_{0}=1$,\quad$T_{1}=x$,\quad$T_{2}=2x^{2}-1$,\quad$(T_{3}=4x^{3}-3x)$.
\solution
Substitute $x=\cos t$. Since $dx=-\sin t\,dt$,
$\sqrt{1-x^2}=\sin t$ on $0\le t\le\pi$, and $x=-1,1$ correspond to
$t=\pi,0$ (the minus sign of $dx$ is absorbed by interchanging these
limits),
\[
M_{2j}=\int_{-1}^1\frac{x^{2j}}{\sqrt{1-x^2}}dx
=\int_0^\pi\cos^{2j}t\,dt,
\quad M_{2j}=\frac{2j-1}{2j}M_{2j-2}.
\]
To see every term of the reduction, integrate
$\cos^{2j-1}t\,d(\sin t)$ by parts:
\[
M_{2j}=[\cos^{2j-1}t\sin t]_0^\pi
 +(2j-1)\int_0^\pi\sin^2t\cos^{2j-2}t\,dt
 =(2j-1)(M_{2j-2}-M_{2j}).
\]
The boundary term vanishes because $\sin0=\sin\pi=0$, and the last step
uses $\sin^2t=1-\cos^2t$. Moving the last term to the left gives
$2jM_{2j}=(2j-1)M_{2j-2}$, the stated recurrence. It yields
$M_0=\pi$, $M_2=\pi/2$, $M_4=3\pi/8$, and $M_6=5\pi/16$;
all odd moments vanish. Gram--Schmidt therefore gives
\[
q_0=1,\qquad q_1=x,\qquad
q_2=x^2-\frac{M_2}{M_0}=x^2-\frac12.
\]
Their squared norms are
\[
\|q_0\|^2=\pi,\quad \|q_1\|^2=M_2=\frac\pi2,\quad
\|q_2\|^2=M_4-M_2+\frac14M_0
=\frac{3\pi}8-\frac\pi2+\frac\pi4=\frac\pi8.
\]
Thus $q_0$ and $q_1$ already have the requested normalization, and $q_2$
needs multiplication by $2$:
\[
\boxed{T_0=1,\qquad T_1=x,\qquad T_2=2x^2-1.}
\]
To check the additional answer $T_3$, only the projection of $x^3$ onto
$x$ survives parity; its coefficient is
$\langle x|x^3\rangle/\langle x|x\rangle=M_4/M_2$:
\[
q_3=x^3-\frac{M_4}{M_2}x=x^3-\frac34x,
\]
\[
\|q_3\|^2=M_6-\frac32M_4+\frac9{16}M_2
=\frac{5\pi}{16}-\frac{9\pi}{16}+\frac{9\pi}{32}
=\frac\pi{32}.
\]
Multiplying $q_3$ by $4$ gives $T_3=4x^3-3x$ and norm squared $\pi/2$.

\Needspace{5\baselineskip}
\par\bigskip
\noindent\textbf{5.2.6}\quad Use the Gram-Schmidt orthogonalization scheme to
construct the first three Chebyshev polynomials (type II):
\[
u_{n}(x)=x^{n},\quad n=0,1,2,\ldots,\quad-1\le x\le1,\quad
w(x)=(1-x^{2})^{+1/2}.
\]
Take the normalization to be
\[
\int_{-1}^{1}U_{m}(x)U_{n}(x)w(x)\,dx=\delta_{mn}\frac{\pi}{2}.
\]
\emph{Hint.}
\[
\int_{-1}^{1}(1-x^{2})^{1/2}x^{2n}\,dx
=\frac{\pi}{2}\times\frac{1\cdot3\cdot5\cdots(2n-1)}
{4\cdot6\cdot8\cdots(2n+2)},\quad n=1,2,3,\ldots
\]
\[
=\frac{\pi}{2},\quad n=0.
\]
\par\nobreak\hfill\emph{ANS.}\enspace
$U_{0}=1$,\quad$U_{1}=2x$,\quad$U_{2}=4x^{2}-1$.
\solution
The hint gives the even moments
\[
M_0=\frac\pi2,\qquad M_2=\frac\pi2\cdot\frac14=\frac\pi8,\qquad
M_4=\frac\pi2\cdot\frac{1\cdot3}{4\cdot6}=\frac\pi{16},
\]
and symmetry makes the odd moments zero. The monic Gram--Schmidt
polynomials are
\[
q_0=1,\qquad q_1=x,\qquad
q_2=x^2-\frac{\langle1|x^2\rangle}{\langle1|1\rangle}
=x^2-\frac14.
\]
There is no projection onto $x$ in the last line because
$\langle x|x^2\rangle=M_3=0$. The squared norms are, with
$q_2^2=x^4-\frac12x^2+\frac1{16}$,
\[
\|q_0\|^2=\frac\pi2,\qquad \|q_1\|^2=M_2=\frac\pi8,
\]
\[
\|q_2\|^2=M_4-\frac12M_2+\frac1{16}M_0
=\frac\pi{16}-\frac\pi{16}+\frac\pi{32}=\frac\pi{32}.
\]
Scale these by $1$, $2$, and $4$ so that every squared norm is $\pi/2$:
\[
\boxed{U_0=1,\qquad U_1=2x,\qquad U_2=4x^2-1.}
\]
The remaining even cross integral is
$\langle U_0|U_2\rangle=4M_2-M_0=0$; the other cross integrals are odd.

\Needspace{5\baselineskip}
\par\bigskip
\noindent\textbf{5.2.7}\quad As a modification of \erratum{Exercise 5.2.4}, apply the
Gram-Schmidt orthogonalization procedure to the set $u_{n}(x)=x^{n}$,
$n=0,1,2,\ldots$, $0\le x<\infty$. Take $w(x)$ to be $\exp(-x^{2})$. Find the
first two nonvanishing polynomials. Normalize so that the coefficient of the
highest power of $x$ is unity. In \erratum{Exercise 5.2.4}, the interval $(-\infty,
\infty)$ led to the Hermite polynomials. The functions found here are
certainly not the Hermite polynomials.
\par\nobreak\hfill\emph{ANS.}\enspace
$\varphi_{0}=1$,\quad$\varphi_{1}=x-\pi^{-1/2}$.
\solution
\erratum{The textbook (p.~275) twice refers to Exercise 5.2.5, which
constructs the Chebyshev polynomials; both references should be to
Exercise 5.2.4, which uses the Gaussian weight on $(-\infty,\infty)$ and
gives the Hermite polynomials.}

On the half-line the moments required for the first two polynomials are
as follows; $M_0$ is half the full-line value $\sqrt\pi$ because the
integrand is even:
\[
M_0=\int_0^\infty e^{-x^2}dx=\frac{\sqrt\pi}{2},\qquad
M_1=\int_0^\infty xe^{-x^2}dx
=\left[-\frac12e^{-x^2}\right]_0^\infty=\frac12.
\]
The monic constant is $\varphi_0=1$, and subtraction of its projection
from $x$ gives
\[
\boxed{\varphi_1=x-\frac{\langle1|x\rangle}{\langle1|1\rangle}
=x-\frac1{\sqrt\pi}.}
\]
Indeed, $\langle\varphi_0|\varphi_1\rangle
=1/2-(1/\sqrt\pi)(\sqrt\pi/2)=0$.
For an independent norm check, $M_2=\sqrt\pi/4$ (half the full-line
value of Exercise 5.2.4), so
\[
\|\varphi_1\|^2=M_2-\frac2{\sqrt\pi}M_1+\frac1\pi M_0
=\frac{\sqrt\pi}4-\frac1{\sqrt\pi}+\frac1{2\sqrt\pi}
=\frac{\pi-2}{4\sqrt\pi}>0.
\]
On the full line the first moment is zero; on the half-line it is $1/2$.
This explicitly explains the constant term absent from the odd Hermite
polynomial.

\Needspace{5\baselineskip}
\par\bigskip
\noindent\textbf{5.2.8}\quad Form a set of three orthonormal vectors by the
Gram-Schmidt process using these input vectors in the order given:
\[
\mathbf{c}_{1}=\begin{pmatrix}1\\1\\1\end{pmatrix},\quad
\mathbf{c}_{2}=\begin{pmatrix}1\\1\\2\end{pmatrix},\quad
\mathbf{c}_{3}=\begin{pmatrix}1\\0\\2\end{pmatrix}.
\]
\solution
Normalize the first vector:
\[
\mathbf e_1=\frac{\mathbf c_1}{\|\mathbf c_1\|}
=\frac1{\sqrt3}(1,1,1)^T.
\]
The projection coefficient of $\mathbf c_2$ onto it is
$\mathbf e_1\cdot\mathbf c_2=4/\sqrt3$. Subtract and normalize:
\[
\mathbf q_2=(1,1,2)^T-\frac43(1,1,1)^T
=\frac13(-1,-1,2)^T,\qquad \|\mathbf q_2\|^2=\frac23,
\]
\[
\mathbf e_2=\frac1{\sqrt6}(-1,-1,2)^T.
\]
For the third vector,
\[
\mathbf e_1\cdot\mathbf c_3=\frac3{\sqrt3}=\sqrt3,
\qquad
\mathbf e_2\cdot\mathbf c_3=\frac{-1+4}{\sqrt6}=\frac3{\sqrt6}.
\]
Since $\sqrt3\,\mathbf e_1=(1,1,1)^T$ and
$(3/\sqrt6)\,\mathbf e_2=\frac12(-1,-1,2)^T$, subtracting both
projections gives
\begin{align*}
\mathbf q_3&=(1,0,2)^T-(1,1,1)^T
 -\frac12(-1,-1,2)^T\\
&=(1/2,-1/2,0)^T,\qquad \|\mathbf q_3\|^2=\frac12,
\end{align*}
so the requested set is
\[
\boxed{\mathbf e_1=\frac{(1,1,1)^T}{\sqrt3},\quad
\mathbf e_2=\frac{(-1,-1,2)^T}{\sqrt6},\quad
\mathbf e_3=\frac{(1,-1,0)^T}{\sqrt2}.}
\]
The three cross-dot-product numerators are $-1-1+2=0$, $1-1=0$,
and $-1+1=0$; their squared lengths are respectively $3/3$, $6/6$,
and $2/2$, independently verifying orthonormality.

\Needspace{9\baselineskip}
\subsection*{5.3 Operators}

\Needspace{5\baselineskip}
\par\bigskip
\noindent\textbf{5.3.1}\quad Show (without introducing matrix representations)
that the adjoint of the adjoint of a \erratum{bounded, everywhere-defined} operator restores the original operator,
i.e., that $(\mathcal{A}^{\dagger})^{\dagger}=\mathcal{A}$.
\solution
\erratum{For the identity as an equality of operators, assume the operators
are bounded and everywhere defined, as they automatically are in finite
dimension. For a densely defined unbounded operator, the double adjoint
is its closure when the operator is closable; equality then additionally
requires that it be closed.}

Let $f,g$ be arbitrary vectors. The adjoint definition and conjugate
symmetry of the scalar product give, step by step,
\begin{align*}
\langle f|Ag\rangle
&=\langle A^\dagger f|g\rangle
=\langle g|A^\dagger f\rangle^*\\
&=\langle(A^\dagger)^\dagger g|f\rangle^*
=\langle f|(A^\dagger)^\dagger g\rangle.
\end{align*}
The first and third equalities are the definition
$\langle u|Bv\rangle=\langle B^\dagger u|v\rangle$, applied with
$(B,u,v)=(A,f,g)$ and $(A^\dagger,g,f)$, respectively; the second and
fourth are $\langle u|v\rangle=\langle v|u\rangle^*$.
Hence $\langle f|[A-(A^\dagger)^\dagger]g\rangle=0$ for every $f$.
For fixed $g$, choose $f=[A-(A^\dagger)^\dagger]g$. Its squared norm is
zero, so it is the zero vector. Since $g$ was arbitrary,
\[
\boxed{(A^\dagger)^\dagger=A.}
\]

\Needspace{5\baselineskip}
\par\bigskip
\noindent\textbf{5.3.2}\quad $\mathcal{U}$ and $\mathcal{V}$ are two arbitrary
\erratum{bounded, everywhere-defined} operators. Without introducing matrix representations of these operators, show
that
\[
(\mathcal{U}\mathcal{V})^{\dagger}
=\mathcal{V}^{\dagger}\mathcal{U}^{\dagger}.
\]
Note the resemblance to adjoint matrices.
\solution
\erratum{The unrestricted equality is valid for bounded, everywhere-defined
operators, including all finite matrices. For general unbounded
operators their domains matter; the formal algebra alone does not prove
equality of the adjoint domains.}
For arbitrary $f,g$ in the bounded setting, apply the adjoint definition
twice, first to $U$ acting on the vector $Vg$, then to $V$ with $U^\dagger f$
in the left slot:
\[
\langle f|UVg\rangle
=\langle U^\dagger f|Vg\rangle
=\langle V^\dagger U^\dagger f|g\rangle.
\]
By definition also $\langle f|UVg\rangle=\langle(UV)^\dagger f|g\rangle$.
Subtracting, $\langle[(UV)^\dagger-V^\dagger U^\dagger]f|g\rangle=0$ for
every $g$; the choice $g=[(UV)^\dagger-V^\dagger U^\dagger]f$ shows, as in
the previous exercise, that this vector is zero. Since $f$ is arbitrary,
\[
\boxed{(UV)^\dagger=V^\dagger U^\dagger.}
\]
The order reverses because $U$, the operator applied last, is the first
to be moved across the scalar product.

\Needspace{5\baselineskip}
\par\bigskip
\noindent\textbf{5.3.3}\quad Consider a Hilbert space spanned by the three
functions $\varphi_{1}=x_{1}$, $\varphi_{2}=x_{2}$, $\varphi_{3}=x_{3}$, and a
scalar product defined by $\langle x_{\nu}|x_{\mu}\rangle=\delta_{\nu\mu}$.
\begin{enumerate}
\item[(a)] Form the $3\times3$ matrix of each of the following operators:
  \[
  \mathcal{A}_{1}=\sum_{i=1}^{3}x_{i}\left(\frac{\partial}{\partial x_{i}}
  \right),\quad
  \mathcal{A}_{2}=x_{1}\left(\frac{\partial}{\partial x_{2}}\right)
  -x_{2}\left(\frac{\partial}{\partial x_{1}}\right).
  \]
\item[(b)] Form the column vector representing $\psi=x_{1}-2x_{2}+3x_{3}$.
\item[(c)] Form the matrix equation corresponding to
  $\chi=(\mathcal{A}_{1}-\mathcal{A}_{2})\psi$ and verify that the matrix
  equation reproduces the result obtained by direct application of
  $\mathcal{A}_{1}-\mathcal{A}_{2}$ to $\psi$.
\end{enumerate}
\solution
(a) Because $\langle x_\nu|x_\mu\rangle=\delta_{\nu\mu}$, the matrix
element $A_{ij}=\langle x_i|Ax_j\rangle$ is the coefficient of $x_i$ in
$Ax_j$: the $j$th matrix column is the expansion of $A\varphi_j$ in the
ordered basis $(x_1,x_2,x_3)$. Since
$\partial x_j/\partial x_i=\delta_{ij}$,
\[
A_1x_j=\sum_i x_i\delta_{ij}=x_j.
\]
For the second operator,
\[
A_2x_1=-x_2,\qquad A_2x_2=x_1,\qquad A_2x_3=0.
\]
Reading off the columns gives
\[
\boxed{A_1=\begin{pmatrix}1&0&0\\0&1&0\\0&0&1\end{pmatrix},\qquad
A_2=\begin{pmatrix}0&1&0\\-1&0&0\\0&0&0\end{pmatrix}.}
\]
(b) The coefficient column of $\psi$ is $c=(1,-2,3)^T$.

(c) Subtraction and multiplication give
\[
(A_1-A_2)c
=\begin{pmatrix}1&-1&0\\1&1&0\\0&0&1\end{pmatrix}
 \begin{pmatrix}1\\-2\\3\end{pmatrix}
=\begin{pmatrix}1+2\\1-2\\3\end{pmatrix}
=\begin{pmatrix}3\\-1\\3\end{pmatrix}.
\]
Independently, differentiating the function itself yields
\[
A_1\psi=x_1-2x_2+3x_3,\qquad
A_2\psi=x_1(-2)-x_2(1)=-2x_1-x_2.
\]
Thus
$\chi=A_1\psi-A_2\psi=(x_1-2x_2+3x_3)-(-2x_1-x_2)=3x_1-x_2+3x_3$,
exactly the function represented by the matrix result $(3,-1,3)^T$.

\Needspace{5\baselineskip}
\par\bigskip
\noindent\textbf{5.3.4}\quad
\begin{enumerate}
\item[(a)] Obtain the matrix representation of $\mathcal{A}=x(d/dx)$ in a
  basis of Legendre polynomials, keeping terms through $P_{3}$. Use the
  orthonormal forms of these polynomials as given in \erratum{Example 5.2.1} and the scalar
  product defined there.
\item[(b)] Expand $x^{3}$ in the orthonormal Legendre polynomial basis.
\item[(c)] Verify that $\mathcal{A}x^{3}$ is given correctly by its matrix
  representation.
\end{enumerate}
\solution
\erratum{The textbook says only ``5.2.1.'' The intended reference is
\emph{Example} 5.2.1 (textbook pp.~270--271), which uses $[-1,1]$.
Exercise 5.2.1 instead uses shifted polynomials on $[0,1]$.}
Write the orthonormal basis without confusing it with the conventionally
normalized $P_n$:
\[
p_0=\frac1{\sqrt2},\quad p_1=\sqrt{\frac32}x,\quad
p_2=\sqrt{\frac52}\frac{3x^2-1}{2},\quad
p_3=\sqrt{\frac72}\frac{5x^3-3x}{2}.
\]
(a) Differentiate each polynomial and then re-expand, using
\[
3x^2=(3x^2-1)+1,\qquad 15x^3-3x=3(5x^3-3x)+6x.
\]
This gives
\begin{align*}
Ap_0&=0,\qquad Ap_1=\sqrt{\frac32}x=p_1,\\
Ap_2&=\sqrt{\frac52}3x^2
=2p_2+\sqrt{\frac52}=2p_2+\sqrt5\,p_0,\\
Ap_3&=\sqrt{\frac72}\frac{15x^3-3x}{2}\\
&=3\sqrt{\frac72}\frac{5x^3-3x}{2}
 +3\sqrt{\frac72}x
=3p_3+\sqrt{21}\,p_1.
\end{align*}
The leftover lower-degree terms are re-expressed with $1=\sqrt2\,p_0$ and
$x=\sqrt{2/3}\,p_1$: $\sqrt{5/2}=\sqrt5\,p_0$ and
$3\sqrt{7/2}\,x=3\sqrt{7/3}\,p_1=\sqrt{21}\,p_1$.
All outputs remain in this four-dimensional span. The matrix element
$A_{ij}=\langle p_i|Ap_j\rangle$ is the coefficient of $p_i$ in $Ap_j$, so
the $j$th column lists the coefficients of $Ap_j$. With row and column
order $0,1,2,3$, the matrix is
\[
\boxed{A=\begin{pmatrix}
0&0&\sqrt5&0\\0&1&0&\sqrt{21}\\0&0&2&0\\0&0&0&3
\end{pmatrix}.}
\]
(b) The coefficients of the even polynomials vanish by parity. The two
odd coefficients are
\begin{align*}
\langle p_1|x^3\rangle
&=\sqrt{\frac32}\int_{-1}^1x^4dx
=\sqrt{\frac32}\frac25=\frac{\sqrt6}{5},\\
\langle p_3|x^3\rangle
&=\frac12\sqrt{\frac72}\int_{-1}^1(5x^6-3x^4)\,dx
=\frac12\sqrt{\frac72}\left(5\frac27-3\frac25\right)\\
&=\frac12\sqrt{\frac72}\,\frac8{35}
=\frac4{35}\,\frac{\sqrt{14}}2=\frac{2\sqrt{14}}{35}.
\end{align*}
Therefore
\[
\boxed{x^3=\frac{\sqrt6}{5}p_1+\frac{2\sqrt{14}}{35}p_3},
\qquad c=\begin{pmatrix}0\\\sqrt6/5\\0\\2\sqrt{14}/35\end{pmatrix}.
\]
Substitution of the polynomials checks the expansion directly: since
\[
\frac{\sqrt6}5\sqrt{\frac32}=\frac{\sqrt9}5=\frac35,\qquad
\frac{2\sqrt{14}}{35}\cdot\frac12\sqrt{\frac72}=\frac{\sqrt{49}}{35}
=\frac15,
\]
the sum is $(3/5)x+(1/5)(5x^3-3x)=x^3$.

(c) Matrix multiplication gives
\[
Ac=\begin{pmatrix}
0\\ \sqrt6/5+\sqrt{21}(2\sqrt{14}/35)\\0\\6\sqrt{14}/35
\end{pmatrix}
=\begin{pmatrix}0\\3\sqrt6/5\\0\\6\sqrt{14}/35\end{pmatrix}=3c,
\]
because $\sqrt{21}\sqrt{14}=\sqrt{294}=7\sqrt6$ makes the second entry
$\sqrt6/5+14\sqrt6/35=\sqrt6/5+2\sqrt6/5=3\sqrt6/5$.
This represents $3x^3$, agreeing with the direct derivative
$x(d/dx)x^3=3x^3$.

\Needspace{9\baselineskip}
\subsection*{5.4 Self-Adjoint Operators}

\Needspace{5\baselineskip}
\par\bigskip
\noindent\textbf{5.4.1}\quad
\erratum{Work with bounded, everywhere-defined operators.}
\begin{enumerate}
\item[(a)] $\mathcal{A}$ is a non-Hermitian operator. Show that the operators
  $\mathcal{A}+\mathcal{A}^{\dagger}$ and $i(\mathcal{A}
  -\mathcal{A}^{\dagger})$ are Hermitian.
\item[(b)] Using the preceding result, show that every non-Hermitian operator
  may be written as a linear combination of two Hermitian operators.
\end{enumerate}
\solution
Use the adjoint identities $(\lambda A)^\dagger=\lambda^*A^\dagger$,
$(A+B)^\dagger=A^\dagger+B^\dagger$, and
$(A^\dagger)^\dagger=A$ \erratum{for bounded, everywhere-defined operators. General unbounded
operators require additional domain assumptions, as noted in
Exercises 5.3.1--5.3.2}.

(a) If $B=A+A^\dagger$, then
\[
B^\dagger=A^\dagger+(A^\dagger)^\dagger=A^\dagger+A=B.
\]
If $C=i(A-A^\dagger)$, conjugating the scalar $i$ is essential:
\[
C^\dagger=-i\{A^\dagger-(A^\dagger)^\dagger\}
=-i(A^\dagger-A)=i(A-A^\dagger)=C.
\]
Thus both operators are Hermitian.

(b) Solve these two definitions for $A$: since
$iC=i^2(A-A^\dagger)=-(A-A^\dagger)$,
\[
B-iC=(A+A^\dagger)+(A-A^\dagger)=2A,\qquad\text{so}\qquad
\boxed{A=\frac12B-\frac i2C.}
\]
Equivalently, define the Hermitian real and imaginary parts
\[
H=\frac{A+A^\dagger}{2},\qquad
K=\frac{A-A^\dagger}{2i};\qquad \boxed{A=H+iK}.
\]
Because $1/(2i)=-i/2$, these are $H=B/2$ and $K=-C/2$, real multiples
of the operators of (a), hence Hermitian, and $H+iK=B/2-iC/2$ is the
first boxed formula. Directly, substitution verifies the reconstruction
because the $A^\dagger$ terms cancel and the two $A/2$ terms add. This
decomposition works also if $A$ happens already to be Hermitian, in
which case $K=0$.

\Needspace{5\baselineskip}
\par\bigskip
\noindent\textbf{5.4.2}\quad Prove that the product of two
\erratum{bounded, everywhere-defined} Hermitian operators is Hermitian
if and only if the two operators commute.
\solution
For Hermitian $A$ and $B$, the product-adjoint identity
$(AB)^\dagger=B^\dagger A^\dagger$ (Exercise 5.3.2) gives
\[
(AB)^\dagger=B^\dagger A^\dagger=BA.
\]
If $AB$ is Hermitian, then $AB=(AB)^\dagger=BA$, so the operators
commute. Conversely, if $AB=BA$, then
$(AB)^\dagger=BA=AB$, so the product is Hermitian. Thus
\[
\boxed{AB\text{ is Hermitian}\ \Longleftrightarrow\ [A,B]=0.}
\]
\erratum{This statement has its usual finite-dimensional or bounded-operator
meaning. For unbounded operators one must also establish the domains of
the product and its adjoint; formal commutation on a subdomain alone is
not a proof of self-adjointness.}

\Needspace{5\baselineskip}
\par\bigskip
\noindent\textbf{5.4.3}\quad $\mathcal{A}$ and $\mathcal{B}$ are noncommuting
\erratum{Hermitian, bounded, everywhere-defined} quantum mechanical operators,
and $\mathcal{C}$ is given by the formula
\[
\mathcal{A}\mathcal{B}-\mathcal{B}\mathcal{A}=i\mathcal{C}.
\]
Show that $\mathcal{C}$ is Hermitian.
\erratum{For an unbounded realization, equality of all adjoint domains used
in the argument must also be established; a formal boundary-condition
assertion alone is insufficient.}
\solution
\erratum{The conclusion requires $A$ and $B$ themselves to be Hermitian,
as when they represent quantum observables. Being merely ``quantum
mechanical operators'' or merely noncommuting does not imply this.}
With that assumption and the indicated adjoint/domain conditions, the
product-adjoint identity of Exercise 5.3.2 gives
\begin{align*}
(AB-BA)^\dagger
&=B^\dagger A^\dagger-A^\dagger B^\dagger\\
&=BA-AB=-(AB-BA)=-iC.
\end{align*}
But the adjoint of the right side of $AB-BA=iC$ is
$(iC)^\dagger=i^*C^\dagger=-iC^\dagger$.
Thus $-iC^\dagger=-iC$, and multiplication by $i$ gives
\[
\boxed{C^\dagger=C.}
\]
Equivalently $C=-i[A,B]$ is Hermitian because $[A,B]$ is
anti-Hermitian. Noncommutation only ensures $C\ne0$; it is not needed
for Hermiticity. To check why the added hypothesis matters, take
\[
A=\begin{pmatrix}0&1\\0&0\end{pmatrix},\quad
B=\begin{pmatrix}1&0\\0&0\end{pmatrix}.
\]
Then $[A,B]=-A$, and $C=iA$ is not Hermitian. This counterexample
rules out the unqualified statement.

\Needspace{5\baselineskip}
\par\bigskip
\noindent\textbf{5.4.4}\quad The operator $\mathcal{L}$ is Hermitian. Show
that $\langle L^{2}\rangle\ge0$, meaning that for all
\erratum{$\psi\in D(L^2)$},
$\langle\psi|L^{2}|\psi\rangle\ge0$.
\solution
\erratum{For an unbounded $L$, the expression in the question requires
$\psi\in D(L^2)=\{\psi\in D(L):L\psi\in D(L)\}$, not merely
$\psi\in D(L)$.}
Write $L^2\psi=L\chi$ with $\chi=L\psi$, move the left factor of $L$
across the scalar product by the adjoint definition
$\langle\psi|L\chi\rangle=\langle L^\dagger\psi|\chi\rangle$, and use
$L^\dagger=L$:
\[
\boxed{\langle\psi|L^2|\psi\rangle
=\langle L^\dagger\psi|L\psi\rangle
=\langle L\psi|L\psi\rangle=\|L\psi\|^2\ge0.}
\]
Equality occurs exactly when $L\psi=0$. For a normalized state this is
the stated expectation value. For an unnormalized nonzero state, dividing
by $\langle\psi|\psi\rangle>0$ preserves the inequality.

\Needspace{5\baselineskip}
\par\bigskip
\noindent\textbf{5.4.5}\quad Consider a Hilbert space whose members are
functions defined on the surface of the unit sphere, with a scalar product of
the form
\[
\langle f|g\rangle=\int d\Omega\,f^{*}g,
\]
where $d\Omega$ is the element of solid angle. Note that the total solid angle
of the sphere is $4\pi$. We work here with the three functions
$\varphi_{1}=Cx/r$, $\varphi_{2}=Cy/r$, $\varphi_{3}=Cz/r$, with $C$ assigned
a value that makes the $\varphi_{i}$ normalized.
\begin{enumerate}
\item[(a)] Find $C$, and show that the $\varphi_{i}$ are also mutually
  orthogonal.
\item[(b)] Form the $3\times3$ matrices of the angular momentum operators
  \[
  L_{x}=-i\left(y\frac{\partial}{\partial z}-z\frac{\partial}{\partial y}
  \right),\quad
  L_{y}=-i\left(z\frac{\partial}{\partial x}-x\frac{\partial}{\partial z}
  \right),\quad
  L_{z}=-i\left(x\frac{\partial}{\partial y}-y\frac{\partial}{\partial x}
  \right).
  \]
\item[(c)] Verify that the matrix representations of the components of
  $\mathbf{L}$ satisfy the angular momentum commutator
  $[L_{x},L_{y}]=iL_{z}$.
\end{enumerate}
\solution
(a) Put $n_x=x/r=\sin\theta\cos\phi$,
$n_y=\sin\theta\sin\phi$, $n_z=\cos\theta$, and
$d\Omega=\sin\theta\,d\theta\,d\phi$. In the $\theta$ integrals
substitute $u=\cos\theta$, $du=-\sin\theta\,d\theta$; the limits
$\theta=0,\pi$ become $u=1,-1$, and reversing them absorbs the minus
sign, while $\sin^2\theta=1-u^2$. With
$\int_0^{2\pi}\cos^2\phi\,d\phi=\int_0^{2\pi}\sin^2\phi\,d\phi=\pi$
(their sum is $2\pi$ and their difference is
$\int_0^{2\pi}\cos2\phi\,d\phi=0$),
\begin{align*}
\int n_x^2d\Omega
&=\left(\int_0^{2\pi}\cos^2\phi\,d\phi\right)
  \left(\int_0^\pi\sin^3\theta\,d\theta\right)\\
&=\pi\int_{-1}^1(1-u^2)du=\pi\left(2-\frac23\right)=\frac{4\pi}{3},\\
\int n_y^2d\Omega&=\frac{4\pi}{3}\quad(\cos^2\phi\to\sin^2\phi),\qquad
\int n_z^2d\Omega=2\pi\int_{-1}^1u^2du=\frac{4\pi}{3}.
\end{align*}
The mixed $xy$ integral contains
$\int_0^{2\pi}\sin\phi\cos\phi\,d\phi=0$; the $xz$ and $yz$
integrals contain respectively the zero integrals of $\cos\phi$ and
$\sin\phi$. Normalization requires $C^2(4\pi/3)=1$ for each
$\varphi_i$, and choosing the common constant real and positive gives
\[
\boxed{C=\sqrt{\frac3{4\pi}}},\qquad
\langle\varphi_i|\varphi_j\rangle=\delta_{ij}.
\]
(b) Rotational derivatives annihilate the radial factor. For example,
\[
\left(y\partial_z-z\partial_y\right)\frac1r
=y\frac{-z}{r^3}-z\frac{-y}{r^3}=0.
\]
The same cancellation occurs for $z\partial_x-x\partial_z$ and
$x\partial_y-y\partial_x$. The product rule therefore leaves only
differentiation of $x,y,z$; for example,
$L_x\varphi_2=-iC(y\cdot0-z\cdot1)/r=iCz/r=i\varphi_3$.
The nine actions are
\[
\begin{array}{c|ccc}
 &\varphi_1&\varphi_2&\varphi_3\\\hline
L_x&0&i\varphi_3&-i\varphi_2\\
L_y&-i\varphi_3&0&i\varphi_1\\
L_z&i\varphi_2&-i\varphi_1&0
\end{array}
\]
Because the $\varphi_i$ are orthonormal by (a), the matrix element
$(L_x)_{ij}=\langle\varphi_i|L_x\varphi_j\rangle$ is the coefficient
of $\varphi_i$ in $L_x\varphi_j$. Thus column $j$ of each matrix lists
the coefficients of the table entry under $\varphi_j$ in that
operator's row; e.g.\ $L_x\varphi_2=i\varphi_3$ gives the second
column $(0,0,i)^T$ of $L_x$:
\[
\boxed{L_x=\begin{pmatrix}0&0&0\\0&0&-i\\0&i&0\end{pmatrix},\quad
L_y=\begin{pmatrix}0&0&i\\0&0&0\\-i&0&0\end{pmatrix},\quad
L_z=\begin{pmatrix}0&-i&0\\i&0&0\\0&0&0\end{pmatrix}.}
\]
Conjugate transposition checks that each is Hermitian.

(c) Each product has a single nonzero entry,
$(L_xL_y)_{21}=(L_x)_{23}(L_y)_{31}=(-i)(-i)=-1$ and
$(L_yL_x)_{12}=(L_y)_{13}(L_x)_{32}=i\cdot i=-1$; every other
row--column pairing meets a zero. Hence
\[
L_xL_y=\begin{pmatrix}0&0&0\\-1&0&0\\0&0&0\end{pmatrix},\qquad
L_yL_x=\begin{pmatrix}0&-1&0\\0&0&0\\0&0&0\end{pmatrix}.
\]
Therefore
\[
[L_x,L_y]=\begin{pmatrix}0&1&0\\-1&0&0\\0&0&0\end{pmatrix}
=i\begin{pmatrix}0&-i&0\\i&0&0\\0&0&0\end{pmatrix}=\boxed{iL_z}.
\]

\Needspace{9\baselineskip}
\subsection*{5.5 Unitary Operators}

\Needspace{5\baselineskip}
\par\bigskip
\noindent\textbf{5.5.1}\quad Show that the \erratum{corrected matrix for Example 5.5.1},
\[
\mathsf U=\frac1{\sqrt2}\begin{pmatrix}
-1&\erratum{i}&0&0&0\\1&\erratum{i}&0&0&0\\
0&0&\erratum{-i}&1&0\\0&0&\erratum{i}&1&0\\
0&0&0&0&\sqrt2
\end{pmatrix},
\]
correctly transforms the vector
$f(\theta,\varphi)=3\chi_{1}+2i\chi_{2}-\chi_{3}+\chi_{5}$ to the
$\{\chi'_{i}\}$ basis by
\begin{enumerate}
\item[(a)] (1) Making a column vector $\mathbf{c}$ that represents
  $f(\theta,\varphi)$ in the $\{\chi_{i}\}$ basis, (2) forming
  $\mathbf{c}'=\mathsf{U}\mathbf{c}$, and (3) comparing the expansion
  $\sum_{i}c'_{i}\chi'_{i}(\theta,\varphi)$ with $f(\theta,\varphi)$;
\item[(b)] Verifying that $\mathsf{U}$ is unitary.
\end{enumerate}
\solution
\erratum{The matrix displayed in Example 5.5.1 (p.~290) has four incorrect
imaginary signs for the basis functions defined on p.~289. In fact its
worked integral gives $u_{22}=+i/\sqrt2$, while the displayed matrix has
$-i/\sqrt2$. The matrix must be corrected before part (a) is true.}
We derive the correction directly from those functions, without relying
on the displayed matrix. Euler's identities give
\begin{align*}
\chi'_1&=-\frac{\chi_1+i\chi_2}{\sqrt2},&
\chi'_2&=\frac{\chi_1-i\chi_2}{\sqrt2},\\
\chi'_3&=\frac{\chi_4+i\chi_3}{\sqrt2},&
\chi'_4&=\frac{\chi_4-i\chi_3}{\sqrt2},&
\chi'_5&=\chi_5.
\end{align*}
For example, $e^{\pm i\phi}=\cos\phi\pm i\sin\phi$ yields the first two
identities, and $e^{2i\phi}=\cos2\phi+i\sin2\phi$ together with
$\sin2\phi=2\sin\phi\cos\phi$ yields the third, with exactly
the normalizing constants in the example. Since
$U_{ij}=\langle\chi'_i|\chi_j\rangle$, the coefficients in these
identities must be conjugated to obtain the rows of $U$: if
$\chi'_i=\sum_ka_{ik}\chi_k$, the orthonormality of the $\chi_k$ gives
$U_{ij}=\sum_ka_{ik}^*\langle\chi_k|\chi_j\rangle=a_{ij}^*$. For instance
$U_{12}=(-i/\sqrt2)^*=i/\sqrt2$. Thus
\[
U=\frac1{\sqrt2}\begin{pmatrix}
-1&\erratum{i}&0&0&0\\
1&\erratum{i}&0&0&0\\
0&0&\erratum{-i}&1&0\\
0&0&\erratum{i}&1&0\\
0&0&0&0&\sqrt2
\end{pmatrix}.
\]
(a) The original and transformed coefficient columns are
\[
c=\begin{pmatrix}3\\2i\\-1\\0\\1\end{pmatrix},\qquad
c'=Uc=\frac1{\sqrt2}
\begin{pmatrix}-3+i(2i)\\3+i(2i)\\(-i)(-1)\\i(-1)\\\sqrt2\end{pmatrix}
=\boxed{\begin{pmatrix}-5/\sqrt2\\1/\sqrt2\\i/\sqrt2\\-i/\sqrt2\\1\end{pmatrix}}.
\]
Substituting the basis relations back into the expansion checks it:
\begin{align*}
\sum_jc'_j\chi'_j
&=\frac52(\chi_1+i\chi_2)+\frac12(\chi_1-i\chi_2)\\
&\quad+\frac i2(\chi_4+i\chi_3)-\frac i2(\chi_4-i\chi_3)+\chi_5\\
&=3\chi_1+2i\chi_2-\chi_3+\chi_5=f,
\end{align*}
since the coefficients of $\chi_1$, $\chi_2$, $\chi_3$ collect to
$\frac52+\frac12$, $\frac{5i-i}2$, and $\frac{i^2}2+\frac{i^2}2$, while
the two $\chi_4$ terms cancel.
For comparison, using the matrix as printed gives
$\widetilde c'=(-1,5,-i,i,\sqrt2)^T/\sqrt2$ and reconstructs
$3\chi_1-2i\chi_2+\chi_3+\chi_5$, which is not the requested $f$.
This explicitly verifies the need for the correction.

(b) The corrected $U$ has two $2\times2$ blocks and a final $1$.
Their products with their adjoints are
\[
\frac12\begin{pmatrix}-1&i\\1&i\end{pmatrix}
\begin{pmatrix}-1&1\\-i&-i\end{pmatrix}
=\frac12\begin{pmatrix}2&-1+1\\-1+1&2\end{pmatrix}=I_2,
\]
\[
\frac12\begin{pmatrix}-i&1\\i&1\end{pmatrix}
\begin{pmatrix}i&-i\\1&1\end{pmatrix}
=\frac12\begin{pmatrix}2&-1+1\\-1+1&2\end{pmatrix}=I_2.
\]
The blocks have disjoint support and $(\sqrt2/\sqrt2)^2=1$ for the last
entry, so $UU^\dagger=I_5$; for a square matrix a right inverse is also
a left inverse, hence $U^\dagger U=I_5$ as well.
The printed matrix is also unitary; that test by itself cannot detect a
matrix that represents the wrong basis transformation.

\Needspace{5\baselineskip}
\par\bigskip
\noindent\textbf{5.5.2}\quad
\begin{enumerate}
\item[(a)] Given (in $\mathbb{R}^{3}$) the basis $\varphi_{1}=x$,
  $\varphi_{2}=y$, $\varphi_{3}=z$, consider the basis transformation
  \erratum{$x=z'$, $y=y'$, $z=-x'$ (old coordinate functions written in new
coordinates)}. Find the $3\times3$ matrix $\mathsf{U}$ for
  this transformation.
\item[(b)] This transformation corresponds to a rotation of the coordinate
  axes. Identify the rotation and reconcile your transformation matrix with an
  appropriate matrix $\mathsf{S}(\alpha,\beta,\gamma)$ of the form given in
  Eq.~(3.37).
\item[(c)] Form the column vector $\mathbf{c}$ representing (in the original
  basis) $f=2x-3y+z$, find the result of applying $\mathsf{U}$ to
  $\mathbf{c}$, and show that this is consistent with the basis transformation
  of part (a).
\end{enumerate}
\emph{Note.}~You do not need to be able to form scalar products to handle this
exercise; a knowledge of the linear relationship between the original and
transformed functions is sufficient.
\solution
\erratum{The textbook writes the transformation as $x\to z$, $y\to y$,
$z\to-x$. The arrows require a convention: here, consistent with the
coefficient transformation in Section 5.5, they express the old
coordinate functions in new coordinates, $x=z'$, $y=y'$, $z=-x'$.
Equivalently the new functions are $x'=-z$, $y'=y$, $z'=x$.}

(a) Each column of $U$ contains the coefficients of an old basis function
in the new basis $(x',y',z')$: for $f=\varphi_j$ the old column is
$c=e_j$, and $c'=Uc$ is column $j$ of $U$. The identities
\[
x=0x'+0y'+1z',\qquad y=0x'+1y'+0z',\qquad
z=-1x'+0y'+0z'
\]
therefore give
\[
\boxed{U=\begin{pmatrix}0&0&-1\\0&1&0\\1&0&0\end{pmatrix}.}
\]
(b) Since $x'=-z$ and $z'=x$, the new coordinate axes are
$\hat e'_x=-\hat e_z$, $\hat e'_y=\hat e_y$,
$\hat e'_z=\hat e_x$. They result from a right-handed rotation of the
axes by $\pi/2$ about $+y$: such a rotation by $\beta$ carries
$\hat e_x\to\cos\beta\,\hat e_x-\sin\beta\,\hat e_z$ and
$\hat e_z\to\sin\beta\,\hat e_x+\cos\beta\,\hat e_z$, which for
$\beta=\pi/2$ are $-\hat e_z$ and $\hat e_x$. For the Euler convention
of Eq.~(3.37), setting $\alpha=\gamma=0$ ($\cos\alpha=\cos\gamma=1$,
$\sin\alpha=\sin\gamma=0$) leaves
\[
S(0,\beta,0)=\begin{pmatrix}
\cos\beta&0&-\sin\beta\\0&1&0\\\sin\beta&0&\cos\beta
\end{pmatrix},
\]
whose rows are the components of the three rotated axes just
described. Putting $\beta=\pi/2$ gives $S(0,\pi/2,0)=U$, as required.

(c) The original coefficient column and its transform are
\[
c=\begin{pmatrix}2\\-3\\1\end{pmatrix},\qquad
c'=Uc=\begin{pmatrix}-1\\-3\\2\end{pmatrix}.
\]
The function itself satisfies
\[
f=2x-3y+z=2z'-3y'-x'=-x'-3y'+2z',
\]
which verifies the coefficient column without scalar products.

If instead the arrows are used to \emph{define} the new ordered basis
as $(z,y,-x)$ in the old coordinates, the coefficient transformation for
the same fixed $f$ is $U^{-1}$ and gives $(1,-3,-2)^T$.
Likewise, active substitution $x\mapsto z$, $z\mapsto-x$ in $f$ produces
$-x-3y+2z$ in the old labels. Keeping the primes explicit distinguishes
these interpretations and their opposite rotation senses.

\Needspace{5\baselineskip}
\par\bigskip
\noindent\textbf{5.5.3}\quad Construct the matrix representing the inverse of
the transformation in Exercise 5.5.2, and show that this matrix and the
transformation matrix of that exercise are matrix inverses of each other.
\solution
Using the convention established in Exercise 5.5.2, solve the
transformation $x=z'$, $y=y'$, $z=-x'$ for the primed variables:
$x'=-z$, $y'=y$, $z'=x$. To express a primed expansion in the old basis,
read the coefficients of these three functions as columns (as in
5.5.2(a), with the roles of the two bases exchanged):
\[
\boxed{V=\begin{pmatrix}0&0&1\\0&1&0\\-1&0&0\end{pmatrix}=U^T=U^\dagger.}
\]
Both orders of multiplication give the identity:
\[
UV=\begin{pmatrix}0&0&-1\\0&1&0\\1&0&0\end{pmatrix}
\begin{pmatrix}0&0&1\\0&1&0\\-1&0&0\end{pmatrix}
=\begin{pmatrix}1&0&0\\0&1&0\\0&0&1\end{pmatrix},
\]
\[
VU=\begin{pmatrix}0&0&1\\0&1&0\\-1&0&0\end{pmatrix}
\begin{pmatrix}0&0&-1\\0&1&0\\1&0&0\end{pmatrix}=I_3.
\]
(Each row of $V$ has a single nonzero entry, so the rows of $VU$ are
row 3 of $U$, row 2 of $U$, and minus row 1 of $U$, namely
$(1,0,0)$, $(0,1,0)$, $(0,0,1)$.)
As a vector check, $V(-1,-3,2)^T=(2,-3,1)^T$, recovering the original
coefficients from 5.5.2.

\Needspace{5\baselineskip}
\par\bigskip
\noindent\textbf{5.5.4}\quad \erratum{Consider two successive changes of coefficient columns,
first from a basis $\{\varphi_i\}$ to $\{\varphi'_i\}$ by $\mathsf U$,
then to $\{\chi_i\}$ by $\mathsf V$. The bases need not be orthonormal.
Take real $\theta$ with
$\sin\theta\cos\theta(\cos^2\theta-\sin^2\theta)\ne0$ so both maps are
invertible.} Their matrix representations are
\[
\mathsf{U}=\begin{pmatrix}i\sin\theta&\cos\theta&0\\
-\cos\theta&i\sin\theta&0\\0&0&1\end{pmatrix},\quad
\mathsf{V}=\begin{pmatrix}1&0&0\\0&\cos\theta&i\sin\theta\\
0&\cos\theta&-i\sin\theta\end{pmatrix}.
\]
Given the function $f(x)=3\varphi_{1}(x)-\varphi_{2}(x)-2\varphi_{3}(x)$,
\begin{enumerate}
\item[(a)] By applying $\mathsf{U}$, form the vector representing $f(x)$ in
  the $\{\varphi'_{i}\}$ basis and then by applying $\mathsf{V}$ form the
  vector representing $f(x)$ in the $\{\chi_{i}\}$ basis. Use this result to
  write $f(x)$ as a linear combination of the $\chi_{i}$.
\item[(b)] Form the matrix products $\mathsf{UV}$ and $\mathsf{VU}$ and then
  apply each to the vector representing $f(x)$ in the $\{\varphi_{i}\}$ basis.
  Verify that the results of these applications differ and that only one of
  them gives the result corresponding to part (a).
\end{enumerate}
\solution
Put $s=\sin\theta$, $c=\cos\theta$, and $a=(3,-1,-2)^T$.
First retain the matrices exactly as printed and evaluate the requested
products.

(a) Row-by-column multiplication gives
\[
Ua=\begin{pmatrix}(is)\cdot3+c\cdot(-1)\\(-c)\cdot3+(is)\cdot(-1)\\-2
\end{pmatrix}
=\begin{pmatrix}3is-c\\-3c-is\\-2\end{pmatrix}.
\]
Applying $V$ to this column gives
\begin{align*}
b=V(Ua)
&=\begin{pmatrix}
3is-c\\c(-3c-is)+is(-2)\\c(-3c-is)-is(-2)
\end{pmatrix}\\
&=\boxed{\begin{pmatrix}
-c+3is\\-3c^2-is(c+2)\\-3c^2+is(2-c)
\end{pmatrix}}.
\end{align*}
Where these matrices are invertible and interpreted as successive
coefficient changes, the formal expansion is
\[
f=(-c+3is)\chi_1+[-3c^2-is(c+2)]\chi_2
                 +[-3c^2+is(2-c)]\chi_3.
\]
\erratum{These signs also correct the Instructor's Manual's
multiplication: the second entry of the initial vector is $-1$, so the
first two entries of $Ua$ are $3is-c$ and $-3c-is$.}

(b) Form both matrix products before acting on $a$; each entry has at
most two nonzero terms, e.g.\ $(UV)_{12}=c\cdot c$,
$(UV)_{23}=(is)(is)=-s^2$, $(VU)_{21}=c\,(-c)$,
$(VU)_{32}=c\,(is)$:
\[
UV=\begin{pmatrix}
is&c^2&ics\\-c&ics&-s^2\\0&c&-is
\end{pmatrix},\qquad
VU=\begin{pmatrix}
is&c&0\\-c^2&ics&is\\-c^2&ics&-is
\end{pmatrix}.
\]
Consequently
\[
(UV)a=\begin{pmatrix}
-c^2+is(3-2c)\\-3c-ics+2s^2\\-c+2is
\end{pmatrix},\qquad
(VU)a=\begin{pmatrix}
-c+3is\\-3c^2-is(c+2)\\-3c^2+is(2-c)
\end{pmatrix}=b.
\]
Only $VU$ corresponds to applying $U$ first and $V$ second. In fact,
the two resulting columns differ for every real $\theta$: equality of
their first entries, $-c^2+3is-2isc=-c+3is$, would require
$c(c-1)+2isc=0$; its real part $c(c-1)$ must vanish, so $c=0$ or $c=1$.
For $c=0$, their second entries are $2$ and $-2is$ and differ; for
$c=1$, $s=0$ and their third entries are $-1$ and $-3$ and differ.

\erratum{The textbook's premise that both printed matrices are unitary
is false; this is why the corrected question permits general bases and
excludes singular angles.}
An explicit check, with
\[
U^\dagger=\begin{pmatrix}-is&-c&0\\c&-is&0\\0&0&1\end{pmatrix},\qquad
V^\dagger=\begin{pmatrix}1&0&0\\0&c&c\\0&-is&is\end{pmatrix},
\]
gives (for instance $(U^\dagger U)_{12}=(-is)c+(-c)(is)=-2isc$ and
$(V^\dagger V)_{33}=(-is)(is)+(is)(-is)=2s^2$)
\[
U^\dagger U=\begin{pmatrix}1&-2isc&0\\2isc&1&0\\0&0&1\end{pmatrix},
\qquad V^\dagger V=\begin{pmatrix}1&0&0\\0&2c^2&0\\0&0&2s^2\end{pmatrix}.
\]
Thus $U$ is unitary only when $sc=0$, whereas $V$ is unitary only when
$c^2=s^2=1/2$. These conditions never hold simultaneously, because
$sc=0$ forces $c^2=0$ or $c^2=1$. Moreover, expanding each determinant
along the row that contains its isolated entry $1$,
\[
\det U=(is)(is)-c(-c)=c^2-s^2,\qquad
\det V=c(-is)-(is)c=-2isc.
\]
At their zeros even an invertible change of basis is impossible. Away
from those zeros the calculations above describe general changes of
coefficients, but not successive changes between orthonormal bases.

\erratum{One consistent repair, not uniquely determined by the printed
problem, is to use the following unitary families instead:}
\[
\widehat U=\begin{pmatrix}is&c&0\\\erratum{c}&is&0\\0&0&1\end{pmatrix},\qquad
\widehat V=\begin{pmatrix}1&0&0\\0&c&is\\0&\erratum{is}&\erratum{c}\end{pmatrix}.
\]
Their nontrivial blocks are $isI+c\sigma_1$ and $cI+is\sigma_1$, with
$\sigma_1=\bigl(\begin{smallmatrix}0&1\\1&0\end{smallmatrix}\bigr)$ and
$\sigma_1^2=I$; multiplication by their adjoints cancels the cross
terms:
\begin{align*}
(-isI+c\sigma_1)(isI+c\sigma_1)
&=s^2I-isc\,\sigma_1+isc\,\sigma_1+c^2\sigma_1^2=(s^2+c^2)I=I,\\
(cI-is\sigma_1)(cI+is\sigma_1)
&=c^2I+isc\,\sigma_1-isc\,\sigma_1+s^2\sigma_1^2=I.
\end{align*}
With these repaired data, the full calculation becomes
\[
\widehat Ua=\begin{pmatrix}-c+3is\\3c-is\\-2\end{pmatrix},\qquad
\widehat V\widehat Ua=\begin{pmatrix}
-c+3is\\3c^2-is(c+2)\\s^2-2c+3ics
\end{pmatrix},
\]
\[
\widehat U\widehat V=\begin{pmatrix}is&c^2&ics\\c&ics&-s^2\\0&is&c\end{pmatrix},
\qquad
\widehat V\widehat U=\begin{pmatrix}is&c&0\\c^2&ics&is\\ics&-s^2&c\end{pmatrix},
\]
\[
\widehat U\widehat Va=\begin{pmatrix}
-c^2+is(3-2c)\\3c-ics+2s^2\\-2c-is
\end{pmatrix}.
\]
The coefficients of the repaired expansion in its final basis are the
entries of $\widehat V\widehat Ua$. These repaired matrices fail to
commute in general; for example $c=0$, $s=1$ makes the second entries
of the two results $-2i$ and $2$. At special angles they may commute
(e.g. $\theta=0$), so an assertion of different products for every angle
would require its own qualification.

\Needspace{5\baselineskip}
\par\bigskip
\noindent\textbf{5.5.5}\quad Three functions which are orthogonal with unit
weight on the range $-1\le x\le1$ are $P_{0}=1$, $P_{1}=x$, and
$P_{2}=\frac{3}{2}x^{2}-\frac{1}{2}$. Another set of functions that are
orthogonal and span the same space are $F_{0}=x^{2}$, $F_{1}=x$,
$F_{2}=5x^{2}-3$. Although much of this exercise can be done by inspection,
write down and evaluate all the integrals that lead to the results when they
are obtained in terms of scalar products.
\begin{enumerate}
\item[(a)] Normalize each of the $P_{i}$ and $F_{i}$.
\item[(b)] Find the unitary matrix $\mathsf{U}$ that transforms from the
  normalized $P_{i}$ basis to the normalized $F_{i}$ basis.
\item[(c)] Find the unitary matrix $\mathsf{V}$ that transforms from the
  normalized $F_{i}$ basis to the normalized $P_{i}$ basis.
\item[(d)] Show that $\mathsf{U}$ and $\mathsf{V}$ are unitary, and that
  $\mathsf{V}=\mathsf{U}^{-1}$.
\item[(e)] Expand $f(x)=5x^{2}-3x+1$ in terms of the normalized versions of
  both bases, and verify that the transformation matrix $\mathsf{U}$ converts
  the $P$-basis expansion of $f(x)$ into its $F$-basis expansion.
\end{enumerate}
\solution
Use $M_k=\int_{-1}^1x^kdx$, so $M_k=0$ for odd $k$ and
$M_k=2/(k+1)$ for even $k$: $M_0=2$, $M_2=2/3$, $M_4=2/5$.

(a) The normalization integrals are
\begin{align*}
\|P_0\|^2&=M_0=2,&\|P_1\|^2&=M_2=\frac23,\\
\|P_2\|^2&=\frac14(9M_4-6M_2+M_0)
=\frac14\left(\frac{18}5-4+2\right)=\frac25,\\
\|F_0\|^2&=M_4=\frac25,&\|F_1\|^2&=M_2=\frac23,\\
\|F_2\|^2&=25M_4-30M_2+9M_0=10-20+18=8.
\end{align*}
The normalized functions are therefore
\[
p_0=\frac1{\sqrt2},\quad p_1=\sqrt{\frac32}x,\quad
p_2=\sqrt{\frac52}\frac{3x^2-1}{2},
\]
\[
g_0=\sqrt{\frac52}x^2,\quad g_1=\sqrt{\frac32}x,\quad
g_2=\frac{5x^2-3}{2\sqrt2}.
\]
For completeness, the original off-diagonal integrals are
\[
\langle P_0|P_1\rangle=M_1=0,\quad
\langle P_0|P_2\rangle=(3M_2-M_0)/2=0,\quad
\langle P_1|P_2\rangle=(3M_3-M_1)/2=0,
\]
\[
\langle F_0|F_1\rangle=M_3=0,\quad
\langle F_0|F_2\rangle=5M_4-3M_2=0,\quad
\langle F_1|F_2\rangle=5M_3-3M_1=0.
\]
(b) With the convention $U_{ij}=\langle\varphi'_i|\varphi_j\rangle$ of
Section 5.5, used in Exercise 5.5.1 (new functions $g_i$, old functions
$p_j$), evaluate all
nine overlaps $U_{ij}=\int_{-1}^1g_ip_jdx$:
\begin{align*}
U_{00}&=\frac{\sqrt5}{2}M_2=\frac{\sqrt5}{3},&
U_{01}&=\frac{\sqrt{15}}2M_3=0,\\
U_{02}&=\frac54(3M_4-M_2)=\frac23,&
U_{10}&=\frac{\sqrt3}{2}M_1=0,\\
U_{11}&=\frac32M_2=1,&
U_{12}&=\frac{\sqrt{15}}4(3M_3-M_1)=0,\\
U_{20}&=\frac14(5M_2-3M_0)=-\frac23,&
U_{21}&=\frac{\sqrt3}{4}(5M_3-3M_1)=0,\\
U_{22}&=\frac{\sqrt5}{8}(15M_4-14M_2+3M_0)
=\frac{\sqrt5}{3}.&&
\end{align*}
In $U_{02}$, $U_{20}$, $U_{22}$ the moments evaluate as
\[
\frac54\Bigl(\frac65-\frac23\Bigr)=\frac54\cdot\frac8{15}=\frac23,\qquad
\frac14\Bigl(\frac{10}3-6\Bigr)=-\frac23,\qquad
\frac{\sqrt5}8\Bigl(6-\frac{28}3+6\Bigr)=\frac{\sqrt5}8\cdot\frac83
=\frac{\sqrt5}3;
\]
for $U_{22}$ the constants combine as
$(2\sqrt2)^{-1}\sqrt{5/2}\cdot\frac12=\sqrt5/8$ and
$(5x^2-3)(3x^2-1)=15x^4-14x^2+3$. Thus
\[
\boxed{U=\begin{pmatrix}\sqrt5/3&0&2/3\\0&1&0\\-2/3&0&\sqrt5/3\end{pmatrix}.}
\]
(c) Reversing the real integrands gives
$V_{ij}=\int p_ig_jdx=U_{ji}$; these are the same nine evaluated
integrals with their indices interchanged. Therefore
\[
\boxed{V=U^T=\begin{pmatrix}\sqrt5/3&0&-2/3\\0&1&0\\2/3&0&\sqrt5/3\end{pmatrix}.}
\]
(d) Let $a=\sqrt5/3$, $b=2/3$, so $a^2+b^2=5/9+4/9=1$.
The $0,2$ block of $UV$ is
\[
\begin{pmatrix}a&b\\-b&a\end{pmatrix}
\begin{pmatrix}a&-b\\b&a\end{pmatrix}
=\begin{pmatrix}a^2+b^2&-ab+ab\\-ab+ab&a^2+b^2\end{pmatrix}=I_2.
\]
The middle entry is $1$ and other entries vanish, so $UV=I_3$; the
$0,2$ block of $VU$ is the same computation with $b\to-b$, so also
$VU=I_3$.
Since $V=U^\dagger$, both matrices are unitary and $V=U^{-1}$.

(e) Evaluate the six coefficient integrals, retaining their moment
expressions to check every term:
\begin{align*}
a_0&=\int p_0f\,dx=\frac{5M_2-3M_1+M_0}{\sqrt2}
=\frac{8\sqrt2}{3},\\
a_1&=\int p_1f\,dx=\sqrt{\frac32}(5M_3-3M_2+M_1)=-\sqrt6,\\
a_2&=\int p_2f\,dx
=\frac12\sqrt{\frac52}(15M_4-9M_3-2M_2+3M_1-M_0)
=\frac{2\sqrt{10}}3,\\
b_0&=\int g_0f\,dx=\sqrt{\frac52}(5M_4-3M_3+M_2)
=\frac{4\sqrt{10}}3,\\
b_1&=\int g_1f\,dx=a_1=-\sqrt6\qquad(g_1=p_1),\\
b_2&=\int g_2f\,dx
=\frac{25M_4-15M_3-10M_2+9M_1-3M_0}{2\sqrt2}
=-\frac{2\sqrt2}{3}.
\end{align*}
(All integrals in this part have limits $-1,1$.) The moment
combinations in $a_0,a_1,a_2,b_0,b_2$ are
\[
\frac{10}3+2=\frac{16}3,\qquad -2,\qquad 6-\frac43-2=\frac83,\qquad
2+\frac23=\frac83,\qquad 10-\frac{20}3-6=-\frac83,
\]
where the integrands of $a_2$ and $b_2$ use
\[
(3x^2-1)f=15x^4-9x^3-2x^2+3x-1,\qquad
(5x^2-3)f=25x^4-15x^3-10x^2+9x-3.
\]
With $\sqrt{3/2}=\sqrt6/2$ and $\sqrt{5/2}=\sqrt{10}/2$ these give the
values shown, e.g.
\[
a_0=\frac{16}{3\sqrt2}=\frac{8\sqrt2}3,\qquad
a_2=\frac12\cdot\frac{\sqrt{10}}2\cdot\frac83=\frac{2\sqrt{10}}3,\qquad
b_2=-\frac83\cdot\frac1{2\sqrt2}=-\frac{2\sqrt2}3.
\]
Thus
\[
\boxed{f=\frac{8\sqrt2}{3}p_0-\sqrt6\,p_1+\frac{2\sqrt{10}}3p_2
=\frac{4\sqrt{10}}3g_0-\sqrt6\,g_1-\frac{2\sqrt2}{3}g_2.}
\]
Direct substitution yields, respectively,
$8/3-3x+(5/3)(3x^2-1)$ and
$(20/3)x^2-3x-(5x^2-3)/3$, both equal to $f$.
Finally,
\[
Ua=\begin{pmatrix}
(\sqrt5/3)(8\sqrt2/3)+(2/3)(2\sqrt{10}/3)\\
-\sqrt6\\
(-2/3)(8\sqrt2/3)+(\sqrt5/3)(2\sqrt{10}/3)
\end{pmatrix}
=\begin{pmatrix}4\sqrt{10}/3\\-\sqrt6\\-2\sqrt2/3\end{pmatrix}=b,
\]
because $\sqrt5\sqrt2=\sqrt{10}$ and $\sqrt5\sqrt{10}=5\sqrt2$ make the
first entry $(8\sqrt{10}+4\sqrt{10})/9=12\sqrt{10}/9$ and the third
$(-16\sqrt2+10\sqrt2)/9=-6\sqrt2/9$. This checks the coefficient
transformation independently of reconstruction.

\Needspace{9\baselineskip}
\subsection*{5.6 Transformations of Operators}

\Needspace{5\baselineskip}
\par\bigskip
\noindent\textbf{5.6.1}\quad
\begin{enumerate}
\item[(a)] Using the two spin functions $\varphi_{1}=\alpha$ and
  $\varphi_{2}=\beta$ as an orthonormal basis (so
  $\langle\alpha|\alpha\rangle=\langle\beta|\beta\rangle=1$,
  $\langle\alpha|\beta\rangle=0$), and the relations
  \[
  S_{x}\alpha=\frac{1}{2}\beta,\quad S_{x}\beta=\frac{1}{2}\alpha,\quad
  S_{y}\alpha=\frac{i}{2}\beta,\quad S_{y}\beta=-\frac{i}{2}\alpha,\quad
  S_{z}\alpha=\frac{1}{2}\alpha,\quad S_{z}\beta=-\frac{1}{2}\beta,
  \]
  construct the $2\times2$ matrices of $S_{x}$, $S_{y}$, and $S_{z}$.
\item[(b)] Taking now the basis $\varphi'_{1}=C(\alpha+\beta)$,
  $\varphi'_{2}=C(\alpha-\beta)$:
  \begin{enumerate}
  \item[(i)] Verify that $\varphi'_{1}$ and $\varphi'_{2}$ are orthogonal,
  \item[(ii)] Assign $C$ a value that makes $\varphi'_{1}$ and $\varphi'_{2}$
    normalized,
  \item[(iii)] Find the unitary matrix for the transformation
    $\{\varphi_{i}\}\to\{\varphi'_{i}\}$.
  \end{enumerate}
\item[(c)] Find the matrices of $S_{x}$, $S_{y}$, and $S_{z}$ in the
  $\{\varphi'_{i}\}$ basis.
\end{enumerate}
\solution
(a) In the order $(\alpha,\beta)$, an operator's first column gives its
action on $\alpha$ and its second gives its action on $\beta$.
The six given relations therefore give
\[
\boxed{S_x=\frac12\begin{pmatrix}0&1\\1&0\end{pmatrix},\quad
S_y=\frac12\begin{pmatrix}0&-i\\i&0\end{pmatrix},\quad
S_z=\frac12\begin{pmatrix}1&0\\0&-1\end{pmatrix}.}
\]
For example, $S_y\alpha=(i/2)\beta$ supplies the column $(0,i/2)^T$;
this fixes the imaginary signs without an appeal to a memorized matrix.

(b) Expand the scalar products using orthonormality:
\[
\langle\alpha+\beta|\alpha-\beta\rangle=1-0+0-1=0,
\]
\[
\|\alpha\pm\beta\|^2=1\pm0\pm0+1=2.
\]
Hence $\|\varphi'_{1,2}\|^2=2C^2=1$, and $C=1/\sqrt2$ is a real
positive normalizing choice. The overlap
matrix $U_{ij}=\langle\varphi'_i|\varphi_j\rangle$ has, for example,
$U_{21}=\langle(\alpha-\beta)/\sqrt2|\alpha\rangle=1/\sqrt2$ and
$U_{22}=\langle(\alpha-\beta)/\sqrt2|\beta\rangle=-1/\sqrt2$; altogether
\[
\boxed{U=\frac1{\sqrt2}\begin{pmatrix}1&1\\1&-1\end{pmatrix}},
\qquad U^\dagger=U,\quad U^2=\frac12\begin{pmatrix}2&0\\0&2\end{pmatrix}=I.
\]
Thus $U^\dagger U=U^2=I$: $U$ is unitary.

(c) For $U_{ij}=\langle\varphi'_i|\varphi_j\rangle$ an operator matrix
transforms as $A'=UAU^\dagger$: expanding in the orthonormal $\varphi_k$,
$\varphi'_j=\sum_k\langle\varphi_k|\varphi'_j\rangle\varphi_k
=\sum_kU_{jk}^*\varphi_k=\sum_k(U^\dagger)_{kj}\varphi_k$, and so
$A'_{ij}=\langle\varphi'_i|A\varphi'_j\rangle
=\sum_{k,l}U_{ik}A_{kl}(U^\dagger)_{lj}$, where conjugate linearity in
the first slot turned $(U^\dagger)_{ki}^*$ into $U_{ik}$.
Hence $S'_j=US_jU^\dagger$. Put $H=\sqrt2U$, so that
$US_jU^\dagger=\frac12HS_jH$ because $U^\dagger=U$.
The three unscaled products are
\begin{align*}
H\begin{pmatrix}0&1\\1&0\end{pmatrix}H
&=\begin{pmatrix}1&1\\-1&1\end{pmatrix}
  \begin{pmatrix}1&1\\1&-1\end{pmatrix}
=\begin{pmatrix}2&0\\0&-2\end{pmatrix},\\
H\begin{pmatrix}0&-i\\i&0\end{pmatrix}H
&=\begin{pmatrix}i&-i\\-i&-i\end{pmatrix}
  \begin{pmatrix}1&1\\1&-1\end{pmatrix}
=\begin{pmatrix}0&2i\\-2i&0\end{pmatrix},\\
H\begin{pmatrix}1&0\\0&-1\end{pmatrix}H
&=\begin{pmatrix}1&-1\\1&1\end{pmatrix}
  \begin{pmatrix}1&1\\1&-1\end{pmatrix}
=\begin{pmatrix}0&2\\2&0\end{pmatrix}.
\end{align*}
Each must be multiplied by $1/4$ (two factors $1/\sqrt2$ and the
operator's $1/2$), giving
\[
\boxed{S'_x=\frac12\begin{pmatrix}1&0\\0&-1\end{pmatrix},\quad
S'_y=\frac12\begin{pmatrix}0&i\\-i&0\end{pmatrix},\quad
S'_z=\frac12\begin{pmatrix}0&1\\1&0\end{pmatrix}.}
\]
A direct action check is $S_x(\alpha\pm\beta)=\pm(\alpha\pm\beta)/2$,
which confirms the diagonal $S'_x$ and the order of its eigenvalues.

\Needspace{5\baselineskip}
\par\bigskip
\noindent\textbf{5.6.2}\quad For the basis $\varphi_{1}=Cxe^{-r^{2}}$,
$\varphi_{2}=Cye^{-r^{2}}$, $\varphi_{3}=Cze^{-r^{2}}$, where
$r^{2}=x^{2}+y^{2}+z^{2}$, with the scalar product defined as an unweighted
integral over $\mathbb{R}^{3}$ and with $C$ chosen to make the $\varphi_{i}$
normalized:
\begin{enumerate}
\item[(a)] Find the $3\times3$ matrix of
  $L_{x}=-i\left(y\dfrac{\partial}{\partial z}-z\dfrac{\partial}{\partial y}
  \right)$;
\item[(b)] Using the transformation matrix
  $\mathsf{U}=\begin{pmatrix}1&0&0\\0&1/\sqrt{2}&-i/\sqrt{2}\\
  0&1/\sqrt{2}&i/\sqrt{2}\end{pmatrix}$, find the transformed matrix of
  $L_{x}$;
\item[(c)] Find the new basis functions $\varphi'_{i}$ defined by the
  transformation $\mathsf{U}$, and write explicitly (in terms of $x$, $y$, and
  $z$) the functional forms of $L_{x}\varphi'_{i}$, $i=1,2,3$.
\end{enumerate}
\emph{Hint.}~Use $\int e^{-r^{2}}d^{3}r=\pi^{3/2}$,
$\int x^{2}e^{-r^{2}}d^{3}r=\frac{1}{2}\pi^{3/2}$; the integrals are over
$\mathbb{R}^{3}$.
\solution
First determine the common normalization, choosing $C$ real and positive. Scaling the Gaussian integral
in the hint gives
\[
\int_{\mathbb R^3}e^{-ar^2}d^3r=\pi^{3/2}a^{-3/2},\qquad
\int_{\mathbb R^3}x^2e^{-ar^2}d^3r=\frac{\pi^{3/2}}{2a^{5/2}}.
\]
For the second formula substitute $(X,Y,Z)=\sqrt a(x,y,z)$ in the
hint: $x^2d^3r=X^2d^3R/a^{5/2}$ and $e^{-ar^2}=e^{-R^2}$. The squared
wavefunctions involve $a=2$, not $a=1$:
$\|\varphi_1\|^2=C^2\int x^2e^{-2r^2}d^3r=C^2\pi^{3/2}/(2\cdot2^{5/2})$,
so
\[
1=C^2\frac{\pi^{3/2}}{2^{7/2}},\qquad
\boxed{C=\frac{2^{7/4}}{\pi^{3/4}}}.
\]
By symmetry the same $C$ normalizes $\varphi_2$ and $\varphi_3$. The
mixed integrals vanish by odd parity (e.g.\ $\int xy\,e^{-2r^2}d^3r$
is odd in $x$), so the basis is orthonormal.

(a) Write $E=e^{-r^2}$. Its rotational derivative is zero:
\[
(y\partial_z-z\partial_y)E=y(-2zE)-z(-2yE)=0.
\]
The product rule then gives
\[
L_x(xE)=0,\qquad L_x(yE)=-i(0-z)E=izE,\qquad
L_x(zE)=-i(y-0)E=-iyE.
\]
The common factor $C$ does not change the coefficients, and in the
orthonormal basis column $j$ holds the coefficients of $L_x\varphi_j$, so
\[
\boxed{L_x=\begin{pmatrix}0&0&0\\0&0&-i\\0&i&0\end{pmatrix}.}
\]
(b) The nontrivial block of $U$ has rows $(1,-i)/\sqrt2$ and
$(1,i)/\sqrt2$, whose norms are one and whose mutual Hermitian scalar
product is $(1+(-i)(-i))/2=0$. Thus $U^{-1}=U^\dagger$, and the
transformed matrix is $L'_x=UL_xU^\dagger$ (Exercise 5.6.1(c)).
Explicitly,
\[
UL_x=\begin{pmatrix}0&0&0\\0&1/\sqrt2&-i/\sqrt2\\
0&-1/\sqrt2&-i/\sqrt2\end{pmatrix},\qquad
U^\dagger=\begin{pmatrix}1&0&0\\0&1/\sqrt2&1/\sqrt2\\
0&i/\sqrt2&-i/\sqrt2\end{pmatrix}.
\]
The first row of $UL_x$ vanishes. Row 2 of $UL_x$ times columns 2 and 3
of $U^\dagger$ gives $\frac12+\frac{(-i)(i)}2=1$ and
$\frac12+\frac{(-i)(-i)}2=0$; row 3 gives
$-\frac12+\frac{(-i)(i)}2=0$ and $-\frac12+\frac{(-i)(-i)}2=-1$. Hence
\[
\boxed{L'_x=UL_xU^\dagger=\begin{pmatrix}0&0&0\\0&1&0\\0&0&-1\end{pmatrix}.}
\]
(c) Because $U_{ij}=\langle\varphi'_i|\varphi_j\rangle$, the new ket
functions are $\varphi'_i=\sum_jU_{ij}^*\varphi_j$ (Exercise 5.6.1(c)),
i.e.\ they use the \emph{conjugates} of the rows of $U$:
\[
\boxed{\varphi'_1=CxE,\quad
\varphi'_2=\frac{C}{\sqrt2}(y+iz)E,\quad
\varphi'_3=\frac{C}{\sqrt2}(y-iz)E.}
\]
Applying the differential operator, using the actions in (a), gives
\begin{align*}
L_x\varphi'_1&=0,\\
L_x\varphi'_2&=\frac C{\sqrt2}[iz+i(-iy)]E
=\frac C{\sqrt2}(y+iz)E=\varphi'_2,\\
L_x\varphi'_3&=\frac C{\sqrt2}[iz-i(-iy)]E
=-\frac C{\sqrt2}(y-iz)E=-\varphi'_3.
\end{align*}
These are exactly the three columns of the diagonal matrix in (b).

\Needspace{5\baselineskip}
\par\bigskip
\noindent\textbf{5.6.3}\quad The Gram-Schmidt process for converting an
arbitrary basis $\chi_{\nu}$ into an orthonormal set $\varphi_{\nu}$ is
described in Section 5.2 in a way that introduces coefficients of the form
$-\langle\varphi_{\mu}|\chi_{\nu}\rangle$. For bases consisting of three
functions, convert the formulation so that $\varphi_{\nu}$ is expressed
entirely in terms of the $\chi_{\mu}$, thereby obtaining an expression for the
upper-triangular matrix $\mathsf{T}$ appearing in Eq.~(5.81).
\solution
Use the complex scalar-product convention of the chapter and define
\[
S_{ij}=\langle\chi_i|\chi_j\rangle,\quad S_{ji}=S_{ij}^*,\qquad
D_1=S_{11},\quad D_2=S_{11}S_{22}-S_{12}S_{21},\quad D_3=\det S.
\]
Since the $\chi_i$ are linearly independent,
$v^\dagger Sv=\sum_{i,j}v_i^*S_{ij}v_j=\|\sum_iv_i\chi_i\|^2>0$ for every
nonzero column $v$; so $S$, and likewise its leading $1\times1$ and
$2\times2$ blocks (the Gram matrices of $\chi_1$ and of
$\chi_1,\chi_2$), are positive definite. Their determinants
$D_1,D_2,D_3$ are products of positive eigenvalues, hence positive, and
all square roots below are well defined and positive. Gram--Schmidt
begins with
\[
\varphi_1=\frac{\chi_1}{\sqrt{D_1}},\qquad
\psi_2=\chi_2-\langle\varphi_1|\chi_2\rangle\varphi_1
=\chi_2-\frac{S_{12}}{D_1}\chi_1 .
\]
Expanding the norm, including complex conjugation, gives with
$a=S_{12}/D_1$, $a^*=S_{21}/D_1$,
\begin{align*}
\|\psi_2\|^2&=S_{22}-aS_{21}-a^*S_{12}+a^*aS_{11}\\
&=S_{22}-\frac{S_{12}S_{21}}{D_1}
 -\frac{S_{21}S_{12}}{D_1}
 +\frac{S_{21}S_{12}S_{11}}{D_1^2}\\
&=S_{22}-\frac{S_{12}S_{21}}{D_1}
=\frac{S_{11}S_{22}-S_{12}S_{21}}{D_1}
=\frac{D_2}{D_1},
\end{align*}
because $S_{11}=D_1$ makes the last term cancel one of the two middle
ones. Consequently
\[
\varphi_2=\frac{\psi_2}{\|\psi_2\|}
=\sqrt{\frac{D_1}{D_2}}\Bigl(\chi_2-\frac{S_{12}}{D_1}\chi_1\Bigr)
=\frac{D_1\chi_2-S_{12}\chi_1}{\sqrt{D_1D_2}}.
\]
For the third step define $K=D_1S_{23}-S_{21}S_{13}$. The two projection
coefficients, with the conjugation on the first argument accounted for
($D_1$ is real and $S_{12}^*=S_{21}$), are
\[
\langle\varphi_1|\chi_3\rangle=\frac{S_{13}}{\sqrt{D_1}},\qquad
\langle\varphi_2|\chi_3\rangle
=\frac{D_1S_{23}-S_{12}^*S_{13}}{\sqrt{D_1D_2}}
=\frac K{\sqrt{D_1D_2}}.
\]
Substitute the already-expanded $\varphi_1,\varphi_2$ into
$\psi_3=\chi_3-\langle\varphi_1|\chi_3\rangle\varphi_1
-\langle\varphi_2|\chi_3\rangle\varphi_2$ and collect the coefficient of
each $\chi_i$:
\begin{align*}
\psi_3
&=\chi_3-\frac{S_{13}}{D_1}\chi_1
 -\frac K{D_1D_2}(D_1\chi_2-S_{12}\chi_1)\\
&=\chi_3-\frac K{D_2}\chi_2
 +\frac{-D_2S_{13}+S_{12}K}{D_1D_2}\chi_1.
\end{align*}
The last numerator simplifies without any reality assumption, since
the two terms $\pm S_{12}S_{21}S_{13}$ cancel:
\begin{align*}
-D_2S_{13}+S_{12}K
&=-(D_1S_{22}-S_{12}S_{21})S_{13}
  +S_{12}(D_1S_{23}-S_{21}S_{13})\\
&=D_1(S_{12}S_{23}-S_{22}S_{13}).
\end{align*}
Thus, with $-K=S_{21}S_{13}-D_1S_{23}$ and the factor $D_1$ cancelled,
the unnormalized third vector is
\[
\psi_3=\chi_3
 +\frac{S_{21}S_{13}-D_1S_{23}}{D_2}\chi_2
 +\frac{S_{12}S_{23}-S_{22}S_{13}}{D_2}\chi_1.
\]
Its norm can be evaluated by subtracting the squared orthogonal
projection lengths: $\chi_3=\psi_3+\langle\varphi_1|\chi_3\rangle\varphi_1
+\langle\varphi_2|\chi_3\rangle\varphi_2$ is a sum of three mutually
orthogonal vectors, so by Pythagoras
$\|\chi_3\|^2=\|\psi_3\|^2+|\langle\varphi_1|\chi_3\rangle|^2
+|\langle\varphi_2|\chi_3\rangle|^2$, i.e.
\[
\|\psi_3\|^2
=S_{33}-\frac{S_{31}S_{13}}{D_1}-\frac{K^*K}{D_1D_2}
=\frac{D_3}{D_2}.
\]
To verify the determinant simplification explicitly, multiply by
$D_1D_2$ and use $K^*=D_1S_{32}-S_{12}S_{31}$:
\begin{align*}
D_1D_2\|\psi_3\|^2
&=D_1(D_1S_{22}-S_{12}S_{21})S_{33}
 -(D_1S_{22}-S_{12}S_{21})S_{31}S_{13}\\
&\quad-(D_1S_{32}-S_{12}S_{31})(D_1S_{23}-S_{21}S_{13})\\
&=D_1\{D_1S_{22}S_{33}-S_{12}S_{21}S_{33}-S_{22}S_{31}S_{13}\\
&\hspace{34mm}-D_1S_{32}S_{23}+S_{32}S_{21}S_{13}
                          +S_{12}S_{31}S_{23}\}\\
&=D_1\det S.
\end{align*}
In the second step the two products $S_{12}S_{21}S_{31}S_{13}$ (from the
second and third terms) cancel and every remaining term contains $D_1$;
with $D_1=S_{11}$ the six terms in braces are exactly the six terms of
the expansion
$\det S=S_{11}S_{22}S_{33}+S_{12}S_{23}S_{31}+S_{13}S_{21}S_{32}
-S_{13}S_{22}S_{31}-S_{12}S_{21}S_{33}-S_{11}S_{23}S_{32}$.
Dividing $\psi_3$ by $\|\psi_3\|=\sqrt{D_3/D_2}$ now expresses all three
$\varphi_j$ in the original $\chi_i$. Since the coefficients of
$\varphi_j=\sum_iT_{ij}\chi_i$ form column $j$, the answer is
\[
\boxed{T=\begin{pmatrix}
\dfrac1{\sqrt{D_1}}&-\dfrac{S_{12}}{\sqrt{D_1D_2}}&
\dfrac{S_{12}S_{23}-S_{22}S_{13}}{\sqrt{D_2D_3}}\\[2ex]
0&\sqrt{\dfrac{D_1}{D_2}}&
\dfrac{S_{21}S_{13}-D_1S_{23}}{\sqrt{D_2D_3}}\\[2ex]
0&0&\sqrt{\dfrac{D_2}{D_3}}
\end{pmatrix}.}
\]
Every column has unit $S$-norm and is $S$-orthogonal to the others,
because $(T^\dagger ST)_{\mu\nu}=\sum_{i,j}T_{i\mu}^*S_{ij}T_{j\nu}
=\langle\varphi_\mu|\varphi_\nu\rangle=\delta_{\mu\nu}$; so
$T^\dagger ST=I$. Finally, if $A_{ij}=\langle\chi_i|\mathcal A\chi_j\rangle$,
expanding both arguments gives
\[
\langle\varphi_\mu|\mathcal A\varphi_\nu\rangle
=\sum_{i,j}T_{i\mu}^*A_{ij}T_{j\nu}=(T^\dagger AT)_{\mu\nu},
\]
which also verifies the orientation of $T$ in Eq.~(5.81).

\Needspace{9\baselineskip}
\subsection*{5.7 Invariants}

\Needspace{5\baselineskip}
\par\bigskip
\noindent\textbf{5.7.1}\quad Using the formal properties of unitary
transformations, show that the \erratum{operator commutator $[x,p]=iI$,
on a common domain of the products, is invariant under unitary conjugation
when the domain is carried along by the same unitary map}.
\solution
Let $x'=UxU^\dagger$, $p'=UpU^\dagger$, with
$U^\dagger U=UU^\dagger=I$. Multiply before simplifying:
\begin{align*}
[x',p']
&=(UxU^\dagger)(UpU^\dagger)-(UpU^\dagger)(UxU^\dagger)\\
&=Ux(U^\dagger U)pU^\dagger-Up(U^\dagger U)xU^\dagger\\
&=U(xp-px)U^\dagger=U(iI)U^\dagger
=iUU^\dagger=\boxed{iI}.
\end{align*}
\erratum{The canonical relation $[x,p]=iI$ cannot hold for finite square
matrices: $\operatorname{tr}(xp-px)=0$, whereas
$\operatorname{tr}(iI_N)=iN\ne0$. It is an operator relation in an
infinite-dimensional space. For unbounded position and momentum the
calculation holds on a common domain of the products, carried to the
new domain by $U$; it is not a claim that those products are defined on
every Hilbert-space vector.}
For example, the usual relation holds for
$x\psi=x\psi(x)$ and $p\psi=-i\psi'(x)$ on smooth rapidly decreasing
functions, since
\[
x(-i\psi')-[-i(x\psi)']=-ix\psi'+i(\psi+x\psi')=i\psi.
\]
This explicitly supplies a domain and realization for the formal relation.

\Needspace{5\baselineskip}
\par\bigskip
\noindent\textbf{5.7.2}\quad The Pauli matrices
\[
\sigma_{1}=\begin{pmatrix}0&1\\1&0\end{pmatrix},\quad
\sigma_{2}=\begin{pmatrix}0&-i\\i&0\end{pmatrix},\quad
\sigma_{3}=\begin{pmatrix}1&0\\0&-1\end{pmatrix},
\]
have commutator $[\sigma_{1},\sigma_{2}]=2i\sigma_{3}$. Show that this
relationship continues to be valid if these matrices are transformed by
\[
\mathsf{U}=\begin{pmatrix}\cos\theta&\sin\theta\\
-\sin\theta&\cos\theta\end{pmatrix}.
\]
\solution
Let $c=\cos\theta$, $s=\sin\theta$. The columns of $U$ are real
orthonormal vectors, so $U^\dagger=U^T=U^{-1}$. Direct multiplication
gives
\begin{align*}
\sigma'_1=U\sigma_1U^T
&=\begin{pmatrix}s&c\\c&-s\end{pmatrix}
  \begin{pmatrix}c&-s\\s&c\end{pmatrix}
=\begin{pmatrix}2sc&c^2-s^2\\c^2-s^2&-2sc\end{pmatrix},\\
\sigma'_2=U\sigma_2U^T
&=\begin{pmatrix}is&-ic\\ic&is\end{pmatrix}
  \begin{pmatrix}c&-s\\s&c\end{pmatrix}
=\begin{pmatrix}0&-i\\i&0\end{pmatrix},\\
\sigma'_3=U\sigma_3U^T
&=\begin{pmatrix}c&-s\\-s&-c\end{pmatrix}
  \begin{pmatrix}c&-s\\s&c\end{pmatrix}
=\begin{pmatrix}c^2-s^2&-2sc\\-2sc&-(c^2-s^2)\end{pmatrix}.
\end{align*}
In $\sigma'_2$ the diagonal entries are $isc-ics=0$ and $-isc+isc=0$,
and $c^2+s^2=1$ reduces the off-diagonal ones $-i(s^2+c^2)$ and
$i(c^2+s^2)$ to $-i$ and $i$.
Write $C=\cos2\theta=c^2-s^2$ and $S=\sin2\theta=2sc$. Then
\[
\sigma'_1\sigma'_2
=\begin{pmatrix}S&C\\C&-S\end{pmatrix}
 \begin{pmatrix}0&-i\\i&0\end{pmatrix}
=\begin{pmatrix}iC&-iS\\-iS&-iC\end{pmatrix}=i\sigma'_3,
\]
\[
\sigma'_2\sigma'_1
=\begin{pmatrix}0&-i\\i&0\end{pmatrix}
 \begin{pmatrix}S&C\\C&-S\end{pmatrix}
=\begin{pmatrix}-iC&iS\\iS&iC\end{pmatrix}=-i\sigma'_3.
\]
Subtracting proves
\[
\boxed{[\sigma'_1,\sigma'_2]=2i\sigma'_3.}
\]
This also follows from $[UAU^\dagger,UBU^\dagger]=U[A,B]U^\dagger$
(the computation of Exercise 5.7.1), but the displayed products check
the signs and factors directly.

\Needspace{5\baselineskip}
\par\bigskip
\noindent\textbf{5.7.3}\quad
\begin{enumerate}
\item[(a)] The operator $L_{x}$ is defined as
  \[
  L_{x}=-i\left(y\frac{\partial}{\partial z}-z\frac{\partial}{\partial y}
  \right).
  \]
  Verify that the basis $\varphi_{1}=Cxe^{-r^{2}}$,
  $\varphi_{2}=Cye^{-r^{2}}$, $\varphi_{3}=Cze^{-r^{2}}$, where
  $r^{2}=x^{2}+y^{2}+z^{2}$, forms a closed set under the operation of
  $L_{x}$, meaning that when $L_{x}$ is applied to any member of this basis
  the result is a function within the basis space, and construct the
  $3\times3$ matrix of $L_{x}$ in this basis from the result of the
  application of $L_{x}$ to each basis function.
\item[(b)] Verify that $L_{x}(x+iy)e^{-r^{2}}=-ze^{-r^{2}}$, and note that
  this result, using the $\{\varphi_{i}\}$ basis, can be written
  $L_{x}(\varphi_{1}+i\varphi_{2})=-\varphi_{3}$.
\item[(c)] Express the equation of part (b) in matrix form, and write the
  matrix equation that results when each of the quantities is transformed
  using the transformation matrix
  \[
  \mathsf{U}=\begin{pmatrix}1&0&0\\0&1/\sqrt{2}&-i/\sqrt{2}\\
  0&1/\sqrt{2}&i/\sqrt{2}\end{pmatrix}.
  \]
\item[(d)] Regarding the transformation $\mathsf{U}$ as producing a new basis
  $\{\varphi'_{i}\}$, find the explicit form (in $x,y,z$) of the
  $\varphi'_{i}$.
\item[(e)] Using the operator form of $L_{x}$ and the explicit forms of the
  $\varphi'_{i}$, verify the validity of the transformed equation found in
  part (c).
\end{enumerate}
\emph{Hint.}~The results of Exercise 5.6.2 may be useful.
\solution
Put $E=e^{-r^2}$ and use the same positive normalization
$C=2^{7/4}/\pi^{3/4}$ as in Exercise 5.6.2.

(a) Since $(y\partial_z-z\partial_y)E=0$, the product rule gives, as in
Exercise 5.6.2(a),
\[
L_x\varphi_1=0,\qquad L_x\varphi_2=i\varphi_3,\qquad
L_x\varphi_3=-i\varphi_2.
\]
Every result belongs to the span, so the set is closed under $L_x$.
Column $j$ of the matrix holds the coefficients of $L_x\varphi_j$, so
\[
L=\begin{pmatrix}0&0&0\\0&0&-i\\0&i&0\end{pmatrix}.
\]
(b) Applying the same actions without the factor $C$ gives
\[
L_x[(x+iy)E]=L_x(xE)+iL_x(yE)=0+i(izE)=-zE.
\]
Multiplying by $C$ produces
$L_x(\varphi_1+i\varphi_2)=-\varphi_3$.

(c) The input and output coefficient columns are $a=(1,i,0)^T$ and
$b=(0,0,-1)^T$, so the original equation is
\[
La=\begin{pmatrix}0&0&0\\0&0&-i\\0&i&0\end{pmatrix}
\begin{pmatrix}1\\i\\0\end{pmatrix}
=\begin{pmatrix}0\\0\\i^2\end{pmatrix}=b.
\]
Transform both vectors and the operator, not just the vectors:
\[
a'=Ua=\begin{pmatrix}1\\i/\sqrt2\\i/\sqrt2\end{pmatrix},\quad
b'=Ub=\begin{pmatrix}0\\i/\sqrt2\\-i/\sqrt2\end{pmatrix},\quad
L'=ULU^\dagger=\begin{pmatrix}0&0&0\\0&1&0\\0&0&-1\end{pmatrix},
\]
with $L'$ taken from Exercise 5.6.2(b). For instance, the second entry
of $a'$ is $\frac1{\sqrt2}\cdot i+\bigl(-\frac i{\sqrt2}\bigr)\cdot0$ and
the third entry of $b'$ is $\frac i{\sqrt2}\cdot(-1)$. The equation
survives the transformation because $U^\dagger U=I$:
$L'a'=ULU^\dagger Ua=U(La)=Ub=b'$.
The transformed equation is therefore
\[
\boxed{\begin{pmatrix}0&0&0\\0&1&0\\0&0&-1\end{pmatrix}
\begin{pmatrix}1\\i/\sqrt2\\i/\sqrt2\end{pmatrix}
=\begin{pmatrix}0\\i/\sqrt2\\-i/\sqrt2\end{pmatrix}.}
\]
(d) Conjugating the rows of $U$ to form the new ket functions,
$\varphi'_i=\sum_jU_{ij}^*\varphi_j$ as in Exercise 5.6.2(c), yields
\[
\varphi'_1=CxE,\qquad
\varphi'_2=\frac C{\sqrt2}(y+iz)E,\qquad
\varphi'_3=\frac C{\sqrt2}(y-iz)E.
\]
(e) Direct differentiation, including the cancellation of the radial
terms as in (a), gives
\[
L_x\varphi'_1=0,\qquad
L_x\varphi'_2=\frac C{\sqrt2}[iz+i(-iy)]E=\varphi'_2,
\]
\[
L_x\varphi'_3=\frac C{\sqrt2}[iz-i(-iy)]E=-\varphi'_3.
\]
The transformed matrix equation thus means exactly
\[
L_x\left(\varphi'_1+\frac i{\sqrt2}\varphi'_2
                    +\frac i{\sqrt2}\varphi'_3\right)
=\frac i{\sqrt2}\varphi'_2-\frac i{\sqrt2}\varphi'_3.
\]
By the three actions just found, the left side is
$0+\frac i{\sqrt2}\varphi'_2+\frac i{\sqrt2}(-\varphi'_3)$, which is the
right side; this verifies the transformed equation. To see that it is
the equation of (b), substitute the explicit forms, using
$\frac i{\sqrt2}\cdot\frac C{\sqrt2}=\frac{iC}2$. The input is
\[
CxE+\frac{iC}{2}[(y+iz)+(y-iz)]E=CxE+\frac{iC}2(2y)E=C(x+iy)E,
\]
and the output is
\[
\frac{iC}{2}[(y+iz)-(y-iz)]E=\frac{iC}2(2iz)E=-CzE.
\]
This is precisely the functional equation from (b), verifying the
invariance using both matrices and differential operators.
